\section{AIME I 2015}
        \begin{enumerate}
        	\item (\textit{AIME I 2015}) The expressions $A$ = $1 \times 2 + 3 \times 4 + 5 \times 6 + \cdots + 37 \times 38 + 39$ and $B$ = $1 + 2 \times 3 + 4 \times 5 + \cdots + 36 \times 37 + 38 \times 39$ are obtained by writing multiplication and addition operators in an alternating pattern between successive integers.  Find the positive difference between integers $A$ and $B$.


 



 
 \textbf{Answer: }722



\textbf{Solution 1}

We have 
 \[|A-B|=|1+3(4-2)+5(6-4)+ \cdots + 37(38-36)-39(1-38)|\]
 \[\implies |2(1+3+5+7+ \cdots +37)-1-39(37)|\]
 \[\implies |361(2)-1-39(37)|=|722-1-1443|=|-722|\implies \boxed{722}\]

 





\textbf{Solution 2}

We see that

 $A=(1\times 2)+(3\times 4)+(5\times 6)+\cdots +(35\times 36)+(37\times 38)+39$

 and

 $B=1+(2\times 3)+(4\times 5)+(6\times 7)+\cdots +(36\times 37)+(38\times 39)$.

 Therefore,

 $B-A=-38+(2\times 2)+(2\times 4)+(2\times 6)+\cdots +(2\times 36)+(2\times 38)$

 $=-38+4\times (1+2+3+\cdots+19)$

 $=-38+4\times\frac{20\cdot 19}{2}=-38+760=\boxed{722}.$

 





\textbf{Solution 3 (slower solution)}

For those that aren't shrewd enough to recognize the above, we may use Newton's Little Formula to semi-bash the equations.

 We write down the pairs of numbers after multiplication and solve each layer:

 

 \[2, 12, 30, 56, 90...(39)\]

 
 \[10, 18, 26, 34...\]

 
 \[8, 8, 8...\]

 and

 
 \[(1) 6, 20, 42, 72...\]
 \[14, 22, 30...\]

 
 \[8, 8, 8...\]

 
Then we use Newton's Little Formula for the sum of $n$ terms in a sequence.

 Notice that there are $19$ terms in each sequence, plus the tails of $39$ and $1$ on the first and second equations, respectively.

 
So,

 

 \[2\binom{19}{1}+10\binom{19}{2}+8\binom{19}{3}+1\]

 

 \[6\binom{19}{1}+14\binom{19}{2}+8\binom{19}{3}+39\]

 
Subtracting $A$ from $B$ gives:

 

 \[4\binom{19}{1}+4\binom{19}{2}-38\]

 
Which unsurprisingly gives us $\boxed{722}.$

 
-jackshi2006

 





\textbf{Solution 4 (Cheezy Peezy!)}

We can stockpile A onto B, like this, so an expression like $(1 \times 2) + (3 \times 4) + (5 \times 6) + ... + (37\times 38) + 39 - [(2 \times 3) + (4 \times 5) + ... + (38 \times 39) +1]$ is formed. This can be achieved by rearranging the one in this expression to be under the 39, so that there are neat even multpiles. 

 Next, using paper, one can stock up the expressions like the image shown below. We notice that we have stacked up a lot of even multiples, such as multpiles of 2, 4, 6, 8, et cetera, going on to 38. 

 Now, notice in the first "column", of $(1 \times 2) - (2 \times 3)$, the bottom number multiplying 2 is always two greater than the top number multiplying 2. This goes for all of the even multpiles, up to 38. This means that when the bottom expression is subtracted from the top expression, we are left with a lot of $-2$s being multiplied by even numbers. It cascades all the way to 38, meaning that the value of the chunk (excluding the 39-1) adds up to $-2\times 2 + -2 \times 4 + -2 \times 6 + ... + -2 \times 38$. 

 By the distributive property, the number is also equivalent to $-2(2+4+6+...+38)$. A two can also be factored out of this expression from the parentheses, giving $-2(2(1+2+...+19))$. Finally, because $-2 \times 2 = 4$, as well as the fact that the sum of the first $n$ positive integers is $\frac{(n)(n+1}{2}$, we get the expression $-4(\frac{(19)(20)}{2}) = -4 \times 380 \div 2 = -4 \times 190 = -760.$ 

 However, there is also another thing at the end; the $39 - 1$, using the 1 that was moved earlier! This means we need to add 38 to the number. Because we are in the negatives, adding 38 gives a lower absolute value. Using arithmetic, $-760 + 38 = -722$, but because the answer is the absolute value, the answer is $\boxed{722}$.

 





\textbf{Solution 5 (Summation Notation)}

We have\begin{align*}A&=39+\sum_{k=1}^{19}(2k-1)2k\\&=39+\sum_{k=1}^{19}4k^2-2k\\\end{align*}
\begin{align*}B&=1+\sum_{k=1}^{19}2k(2k+1)\\&=1+\sum_{k=1}^{19}4k^2+2k\\\end{align*}
Therefore,
 \[B-A=\left(\sum_{k=1}^{19}(4k^2+2k)-(4k^2-2k)\right)-38=\left(\sum_{k=1}^{19}4k\right)-38=760-38=\boxed{722}\]

 




	\item (\textit{AIME I 2015}) The nine delegates to the Economic Cooperation Conference include $2$ officials from Mexico, $3$ officials from Canada, and $4$ officials from the United States. During the opening session, three of the delegates fall asleep. Assuming that the three sleepers were determined randomly, the probability that exactly two of the sleepers are from the same country is $\frac{m}{n}$, where $m$ and $n$ are relatively prime positive integers. Find $m+n$.


 



 
 \textbf{Answer: }139



\textbf{Solution 1 (Principle of Inclusion-Exclusion)}

One of the best ways to solve this problem is to use PIE, or the Principle of Inclusion and Exclusion.

 To start off, we know that the denominator, or the total ways to pick $3$ officials, is $\binom{9}{3} = 84$. Now, we find the number of ways that at least $2$ officials are from the same country and then subtract the number of ways all $3$ officials are from the same country. To start with at least $2$ officials, we know:

 
                \begin{itemize}
                    \item There are $7$ different ways to pick $3$ delegates such that $2$ are from Mexico, simply because there are $9-2=7$ "extra" delegates to choose to be the third sleeper once both from Mexico are sleeping.\item There are $3\times7=21$ ways to pick from Canada, as we choose $2$ of the $3$ Canadians ($\binom{3}{2} = 3$) and then there are $7$ other options for the third sleeper.\item Lastly, there are $6\times7=42$ ways to choose for the United States. We can choose two American officials with $\binom{4}{2} = 6$. Then, there are $7$ options for the third sleeper.
                \end{itemize}


                Now, we want to find the number of ways to have three sleepers from the same country.

 
                \begin{itemize}
                    \item There are no ways for the $3$ sleepers to be from Mexico because there are only $2$ Mexican officials. Hence, we get $0$ ways.\item There is only $1$ way to pick all $3$ from Canada because there are exactly $3$ Canadian officials. We now consider the number of times we originally counted this, which after inspection, is $3$, so we have $1 * 3 = 3$.\item Lastly, there are $4$ ways to choose all $3$ officials from the United States ($\binom{4}{3} = 4$). Once again, we counted this $3$ times, so we have $4*3 = 12$.
                \end{itemize}


                Thus, the fraction is $\frac{7+21+42-0-3-12}{84} = \frac{55}{84}$, and our answer is $55+84=\boxed{139}$.

 Note: Similar to Solution 2, we could also have grouped each of the countries in terms of the ways we have for at least $2$ officials to sleep and for all $3$ to sleep. Mexico would have $7 - 0 = 7$ ways for exactly 2 officials to be there, Canada with $21 - 3 = 18$ ways, and the United States with $42 - 12 = 30$ ways.

 Solution by: armang32324

 





\textbf{Solution 2}

The total number of ways to pick $3$ officials from $9$ total is $\binom{9}{3} = 84$, which will be our denominator. Now we can consider the number of ways for exactly two sleepers to be from the same country for each country individually and add them to find our numerator:

 
                \begin{itemize}
                    \item There are $7$ different ways to pick $3$ delegates such that $2$ are from Mexico, simply because there are $9-2=7$ "extra" delegates to choose to be the third sleeper once both from Mexico are sleeping.\item There are $3\times6=18$ ways to pick from Canada, as each Canadian can be left out once and each time one is left out there are $9-3=6$ "extra" people to select one more sleeper from.\item Lastly, there are $6\times5=30$ ways to choose for the United States. It is easy to count $6$ different ways to pick $2$ of the $4$ Americans, and each time you do there are $9-4=5$ officials left over to choose from.
                \end{itemize}


                Thus, the fraction is $\frac{7+18+30}{84} = \frac{55}{84}$. Since this does not reduce, the answer is $55+84=\boxed{139}$.

 





\textbf{Solution 3}

Like in the solution above, there are $84$ ways to pick  $3$ delegates. We can use casework to find the probability that there aren't exactly $2$ sleepers from a county, then subtract from $1$.

 
                \begin{itemize}
                    \item If no country has at least $2$ delegates sleeping, then every country must have $1$ delegate sleeping. There are $2*3*4=24$ ways for this to happen.\item If all $3$ sleeping delegates are from Canada, there are $\binom{3}{3} = 1$ way.\item If all $3$ are from the US, there are $\binom{4}{3} = 4$ ways.
                \end{itemize}


                So, the probability that there are not exactly $2$ sleepers from one country is $\frac{24+1+4}{84} = \frac{29}{84}$, and the probability that exactly $2$ are from the same country is $1- \frac{29}{84} = \frac{55}{84}.$ Our answer is $55+84=\boxed{139}$.

 





\textbf{Solution 4 (Easy Casework)}

Let us take this step-by-step.

 The probability of having exactly 2 Mexican sleepers is $\frac{2}{9} \cdot \frac{1}{8} \cdot 3$.

 The probability of having exactly 2 Canadian sleepers is $\frac{3}{9} \cdot \frac{2}{8} \cdot \frac{6}{7} \cdot 3$.

 The probability of having exactly 2 American sleepers is $\frac{4}{9} \cdot \frac{3}{8} \cdot \frac{5}{7} \cdot 3$.

 Thus, adding these up our total probability is $\frac{55}{84}$. Adding the numerator and denominator gives us our desired $55+84=\boxed{139}.$

 




	\item (\textit{AIME I 2015}) There is a prime number $p$ such that $16p+1$ is the cube of a positive integer.  Find $p$.


 



 
 \textbf{Answer: }307



\textbf{Solution 1}

Let the positive integer mentioned be $a$, so that $a^3 = 16p+1$. Note that $a$ must be odd, because $16p+1$ is odd.

 Rearrange this expression and factor the left side (this factoring can be done using $(a^3-b^3) = (a-b)(a^2+a b+b^2)$ or synthetic divison once it is realized that $a = 1$ is a root):

 
 \begin{align*} a^3-1 &= 16p\\ (a-1)(a^2+a+1) &= 16p\\ \end{align*}


 Because $a$ is odd, $a-1$ is even and $a^2+a+1$ is odd. If $a^2+a+1$ is odd, $a-1$ must be some multiple of $16$. However, for $a-1$ to be any multiple of $16$ other than $16$ would mean $p$ is not a prime. Therefore, $a-1 = 16$ and $a = 17$.

 Then our other factor, $a^2+a+1$, is the prime $p$:

 
 \begin{align*} (a-1)(a^2+a+1) &= 16p\\ (17-1)(17^2+17+1) &=16p\\ p = 289+17+1 &= \boxed{307} \end{align*}


 





\textbf{Solution 2 (Similar to 1)}

Observe that this is the same as $16p+1=n^3$ for some integer $n$.So:

 
 \begin{align*} 16p &= n^3-1\\ 16p &= n^3-1^3\\ 16p &= (n-1)(n^2+n+1)\\ \end{align*}


 Observe that either $p=n-1$ or $p=n^2+n+1$ because $p$ and $16$ share no factors ($p$ can't be $2$).Let $p=n-1$.Then:

 
 \begin{align*} p &= n-1\\ 16 &= n^2+n+1\\ n^2+n &= 15\\ n(n+1) &= 15\\ \end{align*}


 Which is impossible for integer n. So $p=n^2+n+1$ and

 
 \begin{align*} 16 &= n-1\\ n &= 17\\ p &= 17^2+17+1\\ p = 289+17+1 &= \boxed{307}\\ \end{align*}
 - firebolt360

 





\textbf{Solution 3}

Since $16p+1$ is odd, let $16p+1 = (2a+1)^3$. Therefore, $16p+1 = (2a+1)^3 = 8a^3+12a^2+6a+1$. From this, we get $8p=a(4a^2+6a+3)$. We know $p$ is a prime number and it is not an even number. Since $4a^2+6a+3$ is an odd number, we know that $a=8$.

 Therefore, $p=4a^2+6a+3=4*8^2+6*8+3=\boxed{307}$.

 





\textbf{Solution 4}

Let $16p+1=a^3$. Realize that $a$ congruent to $1\mod 4$, so let $a=4n+1$. Expansion, then division by 4, gets $16n^3+12n^2+3n=4p$. Clearly $n=4m$ for some $m$. Substitution and another division by 4 gets $256m^3+48m^4+3m=p$. Since $p$ is prime and there is a factor of $m$ in the LHS, $m=1$. Therefore, $p=\boxed{307}$.

 





\textbf{Solution 5}

Notice that $16p+1$ must be in the form $(a+1)^3 = a^3 + 3a^2 + 3a + 1$. Thus $16p = a^3 + 3a^2 + 3a$, or $16p = a\cdot (a^2 + 3a + 3)$. Since $p$ must be prime, we either have $p = a$ or $a = 16$. Upon further inspection and/or using the quadratic formula, we can deduce $p \neq a$. Thus we have $a = 16$, and $p = 16^2 + 3\cdot 16 + 3 = \boxed{307}$.

 





\textbf{Solution 6}

Notice that the cube 16p+1 is congruent to 1 mod 16. The only cubic numbers that leave a residue of 1 mod 16 are 1 and 15. Case one: The cube is of the form 16k+1-->Plugging this in, and taking note that p is prime and has only 1 factor gives p=307Case two: The cube is of the form 16k+15--> Plugging this in, we quickly realize that this case is invalid, as that implies p is even, and p=2 doesn't work here

 Hence, $p=\boxed{307}$ is our only answer

 pi\_is\_3.141

 





\textbf{Solution 7}

If $16p+1 =k$, we have $k \equiv 1 \mod 16$, so $k^3 \equiv 1 \mod 16$. If $k=1$ we have $p=0$, which is not prime. If $k=17$ we have $16p+1=4913$, or $p=\boxed{307}$

 





\textbf{Solution 8 (Pattern Recognition)}

Notice that:
 \begin{align*}  1^3 &= 0 +1\\ 2^3 &= 1*7+1 \\ 3^3 &= 2*13+1\\ 4^3 &= 3*21+1\\ 5^3 &= 4*31+1\\ 6^3 &= 5*43+1 \end{align*}


 Here, we can see a clear pattern that $n^3=(n-1)p+1$, where $p$ is some positive (not necessarily prime) integer. Hence, the equation $16p+1=a^3$ can interpret as $17^3 = 16p+1$. Solving it, we got $p=307$. After checking all possible divisors, we will find that $307$ is prime. Hence, we got $p=\boxed{307}$.

 





\textbf{Solution 9 (Slightly Different Modular Arithmetic)}

We see that $16p+1=n^3$ for a positive integer $n$. Subtracting $1$, we can turn this equation into a modular congruence, since $n^3-1$ must be a multiple of $16$. 

 Since $n^3-1\equiv0\pmod{16}$, $n^3\equiv1\pmod{16}$. We observe that $n=1$ is a solution to this congruence, which doesn't work. The next, or most obvious number to try is $n=17$. Plugging this in to our original equation, we get 

 $17^3=16p+1$, yielding $p=\boxed{307}$, which is prime.

 
-among us (countmath1)

 





\textbf{Solution 10 (Lengthy Brute Force)}

Recognizing that AIME answers are $0$ through $999$, the numbers whose cube could even be in contention to be equal to $16p + 1$ are $4-25$. The cubes of $1-3$ are all below $17$. We might consider $1$, but that would result in a $p$ of $0$, which is not prime and does not follow our given conditions. We can also see that $16p + 1$ must be an odd number. Hence, we brute force by looking at all cubes of odd numbers in $5-25$ until we get that the cube of $17$, $4913$, works when $p=\boxed{307}$

 Solution by: armang32324

 





\textbf{Solution 11 (Similar to Solution 7, See Last Part)}

What if we say our perfect cube is $x^3$? Then when we solve for $p$, we get:
 \[\]$16p+1=x^3$, so then $x^3\equiv1\mod 16$.
 \[\]Therefore, $x\equiv1\mod 16$. Let's now test values!
 \[\]If $x=1 \Rightarrow p=0$.
 \[\]If $x=17 \Rightarrow 17^3-1=(17-1)(17^2+17+1)=(16)(307)$. 307 is prime! Therefore, $p$ is $\boxed{307}$.
 \[\]~Namish Dour

 




	\item (\textit{AIME I 2015}) Point $B$ lies on line segment $\overline{AC}$ with $AB=16$ and $BC=4$. Points $D$ and $E$ lie on the same side of line $AC$ forming equilateral triangles $\triangle ABD$ and $\triangle BCE$. Let $M$ be the midpoint of $\overline{AE}$, and $N$ be the midpoint of $\overline{CD}$. The area of $\triangle BMN$ is $x$. Find $x^2$.


 



 
 \textbf{Answer: }507



\textbf{Solution 1}

Let $A$ be the origin, so $B=(16,0)$ and $C=(20,0).$ Using equilateral triangle properties tells us that $D=(8,8\sqrt3)$ and $E=(18,2\sqrt3)$ as well. Therefore, $M=(9,\sqrt3)$ and $N=(14,4\sqrt3).$ Applying the Shoelace Theorem to triangle $BMN$ gives

 
 \[x=\dfrac 1 2 |16\sqrt3+36\sqrt3+0-(0+14\sqrt3+64\sqrt3)| =13\sqrt3,\]

 so $x^2=\boxed{507}.$

 





\textbf{Solution 2}

Note that $AB=DB=16$ and $BE=BC=4$. Also, $\angle ABE = \angle DBC = 120^{\circ}$. Thus, $\triangle ABE \cong \triangle DBC$ by SAS.

 From this, it is clear that a $60^{\circ}$ rotation about $B$ will map $\triangle ABE$ to $\triangle DBC$.This rotation also maps $M$ to $N$. Thus, $BM=BN$ and $\angle MBN=60^{\circ}$. Thus, $\triangle BMN$ is equilateral.

 Using the Law of Cosines on $\triangle ABE$,
 \[AE^2 = 16^2 + 4^2 - 2\cdot 16\cdot 4\cdot\left(-\frac{1}{2}\right)\]
 \[AE = 4\sqrt{21}\]Thus, $AM=ME=2\sqrt{21}$.

 Using Stewart's Theorem on $\triangle ABE$,
 \[AM\cdot ME\cdot AE + AE\cdot BM^2 = BE^2\cdot AM + BA^2\cdot ME\]
 \[BM = 2\sqrt{13}\]

 Calculating the area of $\triangle BMN$,
 \[[BMN] = \frac{\sqrt{3}}{4} BM^2\]
 \[[BMN] = 13\sqrt{3}\]Thus, $x=13\sqrt{3}$, so $x^2 = 507$. Our final answer is $\boxed{507}$.

 Admittedly, this is much more tedious than the coordinate solutions.

 I also noticed that there are two more ways of showing that $\triangle BMN$ is equilateral:

 One way is to show that $\triangle ADB$, $\triangle BMN$, and $\triangle ECB$ are related by a spiral similarity centered at $B$.

 The other way is to use the Mean Geometry Theorem. Note that $\triangle BCE$ and $\triangle BDA$ are similar and have the same orientation. Note that $B$ is the weighted average of $B$ and $B$, $M$ is the weighted average of $E$ and $A$, and $N$ is the weighted average of $C$ and $D$. The weights are the same for all three averages. (The weights are actually just $\frac{1}{2}$ and $\frac{1}{2}$, so these are also unweighted averages.) Thus, by the Mean Geometry Theorem, $\triangle BMN$ is similar to both $\triangle BAD$ and $\triangle BEC$, which means that $\triangle BMN$ is equilateral.

 
Note: A much easier way to go about finding $BM$ without having to use Stewart's Theorem is to simply drop the altitudes from M and E to AC, thus hitting AC at points X and Y. Then clearly AEY and AMX are similar with ratio 2. But we know that $AY = 18 \implies AX = 9 \implies BX = 16-9 = 7$. Additionally, $MX = \frac{1}{2} (2\sqrt{3}) = \sqrt{3}$ from similar triangles meaning we can now just do pythagorean theorem on right triangle $MXB$ to get $MB = \sqrt{52}$ - SuperJJ

 





\textbf{Solution 3}


                    \begin{figure}[H]
                        \centering
                        \includegraphics[width=0.5\textwidth]{Solution/AIME_2015_I_4_Solution_Figure_1.png}
                    \end{figure}
                    $AB = BD, BE = BC, \angle ABE = \angle CBD \implies  \triangle ABE \cong \triangle DBC$

 Medians are equal, so $MB = BN, \angle ABM = \angle DBN \implies$$\angle MBN = \angle ABD - \angle ABM + \angle DBN = 60^\circ \implies$

 $\triangle MNB$ is  equilateral triangle.

 The height of $\triangle BCE$ is $2 \sqrt{3},$ distance from $A$ to midpoint $BC$ is $16 + 2 = 18 \implies \frac {AE^2}{4} =\frac{ (16 + 2)^2 +2^2 \cdot 3}{4} = 81 + 3 = 84.$

 $BM$ is the median of $\triangle ABE \implies$$BM^2 = \frac {AB^2}{2} + \frac {BE^2}{2} - \frac {AE^2}{4}=16 \cdot 8 + 4 \cdot 2 - 84 = 52.$ 

 The area of $\triangle BMN$

 
 \[[BMN] = \frac{\sqrt{3}}{4} BM^2 =13 \sqrt{3} \implies \boxed{\textbf{507}}.\]

 \textbf{vladimir.shelomovskii@gmail.com, vvsss}

 





\textbf{Solution 4 (Computational Geometry)}

Take some time understanding that this solution is not a tautology of Solution 1.

 
                    \begin{figure}[H]
                        \centering
                        \includegraphics[width=0.5\textwidth]{Solution/AIME_2015_I_4_Solution_Figure_1.png}
                    \end{figure}
                    Let $A$ be the centre of the origin in the Cartesian plane. On first glance we are not given any direct correspondence of triangle $BMN$ to the rest of the figure, and thus for time efficiency we use computational geometry to solve this problem.

 The line $AD$ has $y-\text{intercept}$ $(0,0)$ and thus its formula is $y=mx$, where $m = \tan(-\frac{\pi}{3}) = -\tan(\frac{\pi}{3}) = -\sqrt{3}$. Thus, point $D$ is located at an $x$ coordinate being the midpoint $AB$ or $8$, and thus $D$ has coordinates $D(8, -8\sqrt{3})$. Line $CE$ has a positive incline passing through the point $C(20,0)$. It follows that the slope of said line is $\tan(\frac{\pi}{6}) = \sqrt{3}$. From there, the $y-\text{intercept}$ $b$ is determined by $-20\sqrt{3} = b$, and thus point $E$ whose $x$ coordinate is 18 satisfies $E(18,-2\sqrt{3})$.

 Through midpoint formulae on both $CD$ and $AE$ points $M$ and $N$ satisfy $M(14, -4\sqrt{3})$, $N(9, -\sqrt{3})$.

 Finally, Shoelace formula on figure $BMN$ yields the desired $\mathcal{X} = 13\sqrt{3}$, thus implying $\mathcal{X}^2 = \boxed{507}$. $\square$

 




	\item (\textit{AIME I 2015}) In a drawer Sandy has $5$ pairs of socks, each pair a different color.  On Monday, Sandy selects two individual socks at random from the $10$ socks in the drawer.  On Tuesday Sandy selects $2$ of the remaining $8$ socks at random, and on Wednesday two of the remaining $6$ socks at random.  The probability that Wednesday is the first day Sandy selects matching socks is $\frac{m}{n}$, where $m$ and $n$ are relatively prime positive integers. Find $m+n$.


 



 
 \textbf{Answer: }341



\textbf{Solution 1}

The order of the days doesn't matter, therefore we can solve backwards. 

 Let the fifth sock be arbitrary; the probability that the sixth sock matches in color is $\dfrac{1}{9}$. 

 Assuming this, then let the first sock be arbitrary; the probability that the second sock does not match is $\dfrac{6}{7}.$

 The only "hard" part is the third and fourth sock. But that is simple casework. If the third sock's color matches the color of one of the first two socks (which occurs with probability $\dfrac{2}{6} = \dfrac{1}{3}$), then the fourth sock can be arbitrary. Otherwise (with probability $\dfrac{2}{3}$), the fourth sock can be chosen with probability $\dfrac{4}{5}$ (5 socks left, 1 sock that can possibly match the third sock's color). The desired probability is thus
 \[\frac{1}{9} \cdot \frac{6}{7} \cdot \left(\dfrac{1}{3} + \dfrac{2}{3} \cdot \dfrac{4}{5}\right) = \frac{26}{315}.\]The sum is therefore $26+315=\boxed{341}.$

 





\textbf{Solution 2}

The key is to count backwards. First, choose the pair which you pick on Wednesday in $5$ ways. Then there are four pairs of socks for you to pick a pair of on Tuesday, and you don't want to pick a pair. Since there are $4$ pairs, the number of ways to do this is $\dbinom{8}{2}-4$. Then, there are two pairs and two nonmatching socks for you to pick from on Monday, a total of $6$ socks. Since you don't want to pick a pair, the number of ways to do this is $\dbinom{6}{2}-2$. Thus the answer is
 \[\dfrac{\left(5\right)\left(\dbinom{8}{2}-4\right)\left(\dbinom{6}{2}-2\right)}{\dbinom{10}{2}\dbinom{8}{2}\dbinom{6}{2}}=\dfrac{26}{315}.\]$26 + 315 = \boxed{341}$.

 





\textbf{Solution 3}

For the first sock, note that to pick two different socks, can complementary count to get the total, $\binom{10}{2}$ minus the number of pairs (5) to get 
 \[\frac{\binom{10}{2} - 5}{\binom{10}{2}}\]The next steps aren't quite as simple, though.WLOG suppose the socks are (a, a, b, b, c, c, d, d, e, e) and that we chose ab on the first day.This leaves abccddee as the remaining socks. We can think of our next choice as how many "pairs" we "break".Note that right now, in abccddee there are 3 pairs, namely {cc, dd, ee}. Our next choice could either leave all of these pairs (if we choose ab), leave 2 of these pairs, (choose say ac to keep dd and ee as pairs), or leave only 1 as a pair, say choosing de to get abccde.
 \[\bold{\text{Case 1: Break zero pairs}}.\] In this case, we want to keep 3 pairs remaining, so our only choice is to choose AB, which would make the remaining ccddee. The probability that we choose AB is $\frac{1}{\binom{8}{2}}$ and the probability that we get a pair after this is $\frac{3}{\binom{6}{2}}$ since we just choose CC, DD, or EE. This case has a probability $\frac{1}{\binom{8}{2}}\frac{3}{\binom{6}{2}}.$
 \[\bold{\text{Case 2: Break one pair}}.\]In this case, we want to choose A and one of {C,C,D,D,E,E} or B and one of {C,C,D,D,E,E}. This yields $\frac{6 + 6}{\binom{8}{2}}$ probability. After this, we end up with 2 pairs, so the total probability of choosing a pair on Wednesday in this case is $\frac{6 + 6}{\binom{8}{2}}\frac{2}{\binom{6}{2}}.$
 \[\bold{\text{case 3: Break 2 pairs}}.\]In this case we'd choose two distinct characters from {C,C,D,D,E,E} which would be $\frac{\binom{6}{2} - 3}{\binom{6}{2}}.$ After this, only one pair remains so the total probability for this case is

 $\frac{\binom{6}{2} - 3}{\binom{6}{2}}    \frac{6}{\binom{8}{2}}\frac{1}{\binom{6}{2}}.$Adding all the cases and multiplying by $\frac{\binom{10}{2} - 5}{\binom{10}{2}}$ yields
 \[\frac{\binom{10}{2} - 5}{\binom{10}{2}}\left(\frac{1}{\binom{8}{2}}\frac{3}{\binom{6}{2}}+\frac{6 + 6}{\binom{8}{2}}\frac{2}{\binom{6}{2}}+ \frac{\binom{6}{2} - 3}{\binom{6}{2}}\frac{6}{\binom{8}{2}}\frac{1}{\binom{6}{2}}\right) = \frac{26}{315} \rightarrow 26 + 315 = \boxed{341}\]

 





\textbf{Solution 4}

This problem is very very similar to \href{https://artofproblemsolving.com/wiki/index.php/2014\_UNCO\_Math\_Contest\_II\_Problems/Problem\_7}{https://artofproblemsolving.com/wiki/index.php/2014\_UNCO\_Math\_Contest\_II\_Problems/Problem\_7}. I will use the results I derived there.We have already established that the probability we get a mismatch on Wednesday GIVEN that there is a mismatch on Monday, Tuesday is $\frac{112}{125}$, meaning the probability for match is $\frac{13}{125}$. Now, we find the probability that Monday,Tuesday are a mismatch and product it with this.Monday: 
 \[1 - \frac{5}{\binom{10}{2}} = \frac{8}{9}\]Tuesday: 
 \[1 - \frac{3}{\binom{8}{2}} = \frac{25}{28}\]
 \[P(\text{ Monday,Tuesday mismatch AND Wednesday match}) = P(\text{ Monday,Tuesday mismatch)} \cdot P(\text{Wednesday match given Monday,Tuesday mismatch})\] 
 \[\implies P = \frac{8}{9} \cdot \frac{25}{28} \cdot \frac{13}{125} = \boxed{\frac{26}{315}}\]

 




	\item (\textit{AIME I 2015}) Point $A,B,C,D,$ and $E$ are equally spaced on a minor arc of a circle. Points $E,F,G,H,I$ and $A$ are equally spaced on a minor arc of a second circle with center $C$ as shown in the figure below. The angle $\angle ABD$ exceeds $\angle AHG$ by $12^\circ$. Find the degree measure of $\angle BAG$.


 
                    \begin{figure}[H]
                        \centering
                        \includegraphics[width=0.5\textwidth]{Problem/AIME_2015_I_6_Problem_Figure_1.png}
                    \end{figure}
                    


 



 
 \textbf{Answer: }058



\textbf{Solution 1}

Let $O$ be the center of the circle with $ABCDE$ on it. 

 Let $x$ be the degree measurement of $\overarc{ED}=\overarc{DC}=\overarc{CB}=\overarc{BA}$ in circle $O$ 

 and $y$ be the degree measurement of $\overarc{EF}=\overarc{FG}=\overarc{GH}=\overarc{HI}=\overarc{IA}$ in circle $C$.

 $\angle ECA$ is, therefore, $5y$ by way of circle $C$ and 
 \[\frac{360-4x}{2}=180-2x\] by way of circle $O$.$\angle ABD$ is $180 - \frac{3x}{2}$ by way of circle $O$, and 
 \[\angle AHG = 180 - \frac{3y}{2}\] by way of circle $C$.

 This means that:

 
 \[180-\frac{3x}{2}=180-\frac{3y}{2}+12\]

 which when simplified yields 
 \[\frac{3x}{2}+12=\frac{3y}{2}\] or 
 \[x+8=y\]Since:
 \[\angle ACE=5y=180-2x\] and 
 \[5x+40=180-2x\]So:
 \[7x=140\Longleftrightarrow x=20\] 
 \[y=28\]$\angle BAG$ is equal to $\angle BAE$ + $\angle EAG$, which equates to $\frac{3x}{2} + y$.Plugging in yields $30+28$, or $\boxed{058}$.

 





\textbf{Solution 2}

Let $m$ be the degree measurement of $\angle GCH$. Since $G,H$ lie on a circle with center $C$, $\angle GHC=\frac{180-m}{2}=90-\frac{m}{2}$. 

 Since $\angle ACH=2 \angle GCH=2m$, $\angle AHC=\frac{180-2m}{2}=90-m$. Adding $\angle GHC$ and $\angle AHC$ gives $\angle AHG=180-\frac{3m}{2}$, and $\angle ABD=\angle AHG+12=192-\frac{3m}{2}$. Since $AE$ is parallel to $BD$, $\angle DBA=180-\angle ABD=\frac{3m}{2}-12=$$\overarc{BE}$. 

 We are given that $A,B,C,D,E$ are evenly distributed on a circle. Hence, 

 $\overarc{ED}=\overarc{DC}=\overarc{CB}=\overarc{BA}$$=\frac{\angle DBA}{3}=\frac{m}{2}-4$

 Here comes the key: Draw a line through $C$ parallel to $AE$, and select a point $X$ to the right of point $C$. 

 $\angle ACX$ = $\overarc{AB}$ + $\overarc{BC}$ = $m-8$.

 Let the midpoint of $\overline{HG}$ be $Y$, then $\angle YCX=\angle ACX+\angle ACY=(m-8)+\frac{5m}{2}=90$. Solving gives $m=28$

 The rest of the solution proceeds as in solution 1, which gives $\boxed{058}$

 


 





\textbf{Solution 3}


                    \begin{figure}[H]
                        \centering
                        \includegraphics[width=0.5\textwidth]{Solution/AIME_2015_I_6_Solution_Figure_1.png}
                    \end{figure}
                    Let $\angle GAH = \varphi \implies \overset{\Large\frown} {GH} = 2\varphi \implies$
 \[\overset{\Large\frown} {EF} = \overset{\Large\frown} {FG} = \overset{\Large\frown} {HI} = \overset{\Large\frown} {IA} = 2\varphi \implies\]
 \[\angle AGH = 2\varphi, \angle ACE = 10 \varphi.\]

 
 \[BD||GH \implies \angle AJB = \angle AGH = 2 \varphi.\]
 \[\triangle AHG: \hspace{10mm}  \angle AHG = \beta = 180^\circ – 3 \varphi.\]$\hspace{10mm} \triangle ABJ:  \hspace{10mm} \angle BAG + \angle ABD = \alpha + \gamma = 180^\circ + 2 \varphi.$

 Let arc $\overset{\Large\frown} {AB} = 2\psi \implies$

 $\angle ACE = \frac {360^\circ – 8 \psi}{2}= 180^\circ – 4 \psi,  \angle ABD = \gamma =\frac {360^\circ – 6 \psi}{2} =180^\circ – 3 \psi.$$\gamma –  \beta = 3(\varphi – \psi) = 12^\circ \implies \psi = \varphi – 4^\circ \implies 10 \varphi = 180^\circ – 4(\varphi – 4^\circ) \implies 14 \varphi = 196^\circ \implies \varphi = 14^\circ.$

 Therefore $\gamma = 180^\circ – 3 \cdot (14^\circ – 4^\circ) = 150^\circ \implies  \alpha = 180^\circ + 2 \cdot 14^\circ – 150^\circ = \boxed{\textbf{058}}.$

 


 


 




	\item (\textit{AIME I 2015}) In the diagram below, $ABCD$ is a square. Point $E$ is the midpoint of $\overline{AD}$. Points $F$ and $G$ lie on $\overline{CE}$, and $H$ and $J$ lie on $\overline{AB}$ and $\overline{BC}$, respectively, so that $FGHJ$ is a square. Points $K$ and $L$ lie on $\overline{GH}$, and $M$ and $N$ lie on $\overline{AD}$ and $\overline{AB}$, respectively, so that $KLMN$ is a square. The area of $KLMN$ is 99. Find the area of $FGHJ$.


 
                    \begin{figure}[H]
                        \centering
                        \includegraphics[width=0.5\textwidth]{Problem/AIME_2015_I_7_Problem_Figure_1.png}
                    \end{figure}
                    


 



 
 \textbf{Answer: }539



\textbf{Solution 1}

Let us find the proportion of the side length of $KLMN$ and $FJGH$. Let the side length of $KLMN=y$ and the side length of $FJGH=x$.

 
                    \begin{figure}[H]
                        \centering
                        \includegraphics[width=0.5\textwidth]{Solution/AIME_2015_I_7_Solution_Figure_1.png}
                    \end{figure}
                    

 
Now, examine $BC$. We know $BC=BJ+JC$, and triangles $\Delta BHJ$ and $\Delta JFC$ are similar to $\Delta EDC$ since they are $1-2-\sqrt{5}$ triangles. Thus, we can rewrite $BC$ in terms of the side length of $FJGH$. 
 \[BJ=\frac{1}{\sqrt{5}}HJ=\frac{x}{\sqrt{5}}=\frac{x\sqrt{5}}{5}, JC=\frac{\sqrt{5}}{2}JF=\frac{x\sqrt{5}}{2}\Rightarrow BC=\frac{7x\sqrt{5}}{10}\]

 Now examine $AB$. We can express this length in terms of $x,y$ since $AB=AN+NH+HB$. By using similar triangles as in the first part, we have 
 \[AB=\frac{1}{\sqrt{5}}y+\frac{\sqrt{5}}{2}y+\frac{2}{\sqrt{5}}x\]
 \[AB=BC\Rightarrow \frac{7y\sqrt{5}}{10}+\frac{2x\sqrt{5}}{5}=\frac{7x\sqrt{5}}{10}\Rightarrow \frac{7y\sqrt{5}}{10}=\frac{3x\sqrt{5}}{10}\Rightarrow 7y=3x\]

 Now, it is trivial to see that $[FJGH]=\left(\frac{x}{y}\right)^2[KLMN]=\left(\frac{7}{3}\right)^2\cdot 99=\boxed{539}.$

 





\textbf{Solution 2}


                    \begin{figure}[H]
                        \centering
                        \includegraphics[width=0.5\textwidth]{Solution/AIME_2015_I_7_Solution_Figure_1.png}
                    \end{figure}
                    

 We begin by denoting the length $ED$ $a$, giving us $DC = 2a$ and $EC = a\sqrt5$. Since angles $\angle DCE$ and $\angle FCJ$ are complementary, we have that $\triangle CDE \sim \triangle JFC$ (and similarly the rest of the triangles are $1-2-\sqrt5$ triangles). We let the sidelength of $FGHJ$ be $b$, giving us: 

 
 \[JC = \sqrt5 \cdot FC = \sqrt5 \cdot FJ/2 = \frac{b\sqrt 5}{2}\]and
 \[BJ = \frac{1}{\sqrt5} \cdot HJ = \frac{b}{\sqrt5}\]

 Since $BC = CJ + BJ$, 

 
 \[2a = \frac{b\sqrt 5}{2} + \frac{b}{\sqrt5}\]

 Solving for $b$ in terms of $a$ yields 
 \[b = \frac{4a\sqrt5}{7}\]

 We now use the given that $[KLMN] = 99$, implying that $KL = LM = MN = NK = 3\sqrt{11}$. We also draw the perpendicular from $E$ to $ML$ and label the point of intersection $P$ as in the diagram at the top

 This gives that 
 \[AM = 2 \cdot AN = 2 \cdot \frac{3\sqrt{11}}{\sqrt5}\]and 
 \[ME = \sqrt5 \cdot MP = \sqrt5 \cdot \frac{EP}{2} = \sqrt5 \cdot \frac{LG}{2} = \sqrt5 \cdot \frac{HG - HK - KL}{2} = \sqrt{5} \cdot \frac{\frac{4a\sqrt5}{7} - \frac{9\sqrt{11}}{2}}{2}\]

 Since $AE$ = $AM + ME$, we get

 
 \[2 \cdot \frac{3\sqrt{11}}{\sqrt5} + \sqrt{5} \cdot \frac{\frac{4a\sqrt5}{7} - \frac{9\sqrt{11}}{2}}{2} = a\]

 
 \[\Rightarrow 12\sqrt{11} + 5(\frac{4a\sqrt5}{7} - \frac{9\sqrt{11}}{2}) = 2\sqrt5a\]

 
 \[\Rightarrow \frac{-21}{2}\sqrt{11} + \frac{20a\sqrt5}{7} = 2\sqrt5a\]

 
 \[\Rightarrow -21\sqrt{11} = 2\sqrt5a\frac{14 - 20}{7}\]

 
 \[\Rightarrow \frac{49\sqrt{11}}{4} = \sqrt5a\]

 
 \[\Rightarrow 7\sqrt{11} = \frac{4a\sqrt{5}}{7}\]

 So our final answer is $(7\sqrt{11})^2 = \boxed{539}$.

 





\textbf{Solution 3}

This is the quickest solution. We see that with a ruler, $KL = \frac{3}{2}$ cm and $HG = \frac{7}{2}$ cm. Thus if $KL$ corresponds with an area of $99$, then $HG$ ($FGHJ$'s area) would correspond with $99*(\frac{7}{3})^2 = \boxed{539}$ - aops5234

 




	\item (\textit{AIME I 2015}) For positive integer $n$, let $s(n)$ denote the sum of the digits of $n$.  Find the smallest positive integer satisfying $s(n) = s(n+864) = 20$.


 



 
 \textbf{Answer: }695



\textbf{Solution 1}

You know whatever $n$ is, it has to have 3 digits, because if it had only two, the maximum of $s(n)$ is 18 and because in AIME all answers are up to three digits. 

 Now let $n=100a_2+10a_1+a_0$

 So first we know, $a_2+a_1+a_0=20$. Okay now we have to split into cases based on which digit gets carried. This meaning, when you add a 3 digit number to 864, we have to know when to carry the digits. Note that if you don't understand any of the steps I take, just try adding any 3-digit number to 864 regularly (using the old-fashioned "put one number over the other" method, not mental calculation), and observe what you do at each step.

 (1)$\textcolor{red}{*}$ None of the digits get carried over to the next space:So this means $a_2<2, a_1<4$ and $a_0<6$. So 

 $s(864+n)=(8+a_2)+(6+a_1)+(4+a_0)=38$So it doesn't work. Now:

 (2) $a_2+8$ is the only one that carries overSo this means $a_2>1, a_1<4$ and $a_0<6$. So

 $s(864+n)=1+(8+a_2-10)+(6+a_1)+(a_0+4)=29$ 

 (3)$\textcolor{red}{*}$ $a_0+4$ is the only one that carries over.So

 $s(864+n)=(8+a_2)+(6+a_1+1)+(4+a_0-10)=29$

 (4)The first and second digit carry over (but not the third)

 $s(864+n)=1+(8+a_2-10+1)+(6+a_1-10)+(4+a_0)=20$

 Aha! This case works but we still have to make sure it's possible for $a_2+a_1+a_0=20$ (We assumed this is true, so we have to find a number that works.) Since only the second and first digit carry over, $a_2>0, a_1>3$ and $a_0<6$. The smallest value we can get with this is 695. Let's see if we can find a smaller one:

 (5)The first and third digit carry over (but not the second)

 $s(864+n)=1+(8+a_2-10)+(7+a_1)+(4+a_0-10)=20$

 The largest value for the middle digit is 2, so the other digits have to be both 9's. So the smallest possible value is 929

 (6) All the digits carry over

 $s(864+n)=1+(9+a_2-10)+(7+a_1-10)+(4+a_0-10)=\text{Way less than 20}$

 
So the answer is $\boxed{695}$ which after a quick test, does indeed work.

 Note: This problem is VERY easy to bash out. I did this when I mocked this test, it gave me the answer in 5 min. Basically you just bash out all of the three digit numbers whose digit sum is 20 and you'll find (quickly) that the answer is $\boxed{695}$

 $\textcolor{red}{*}$ Note to author: Since we have $a_2+a_1+a_0=20$, this implies that $a_2$ must be at least $2$ (since $a_1+a_0 \le 18$), and therefore the first digit must ALWAYS carry. This is presumably the reason where the case where the second and third digits were carried but not the first was omitted, but for this reason cases (1) and (3) are also unnecessary.

 





\textbf{Solution 2}

First, it is easy to verify that $695$ works and that no other numbers beginning with the digit 6 work (i.e. $686, 677, 668, 659$ do not work).

 Suppose by contradiction that there is a smaller valid $n$, where the leading digit of the three-digit number $n$ is 5 or less. (Two-digit $n$ obviously do not work because 9 + 9 < 20.) Clearly $n > 200$ because the smallest three-digit number whose digits sum to 20 is $299$. Also, because the second digit is at most 9, the units digit is at least 6, which means that the addition $N = n + 864$ regroups in the ones place. Then the units digit of $N$ is clearly less than 4. But as $1000 < 200 + 864 < N < 600 + 864 = 1464$, the sum of the thousands digit and the hundredth digit is at most 5. Because the second digit is at most 9, the sum of the digits of $N$ is at most $5 + 9 + 4 < 20$, contradiction. Hence $\boxed{695}$ is the answer.

 \textbf{Modification for Solution 2} During the real test, I immediately noticed that $n$ must be less than 1000 (AIME problem) and that $n$ must be a three-digit number. Therefore, I began casework on the leading digit of $n$. The casework was not intensive (how many ways are there to have digits sum to 20?) and I eventually got 695 as my answer. The rigorous proof that 695 was the smallest came afterwards.

 





\textbf{Solution 3}

First of all, notice that the smallest $n$ with $s(n) = 20$ is $299$. Also, if $s(n + 864) = 20$, $s(n - 136) = 19$ (because subtracting $1000$ from the number removes the $1$ in the thousands place). After checking $s(n - 136)$ for various $n$ with $s(n) = 20$, we see that we need to have a carry when subtracting $136$. To have this, we must either have a $2$ in the tens place or a $5$ in the units place. The minimum $n$ for the former is $929$, and for the latter it is $695$. We check and see that $s(695-136) = s(559) = 19$, so our answer is $\boxed{695}$.

 





\textbf{Solution 4}

Observation (Lemma) : If r is the number of regroups in the addition of n+k, $S(n+k) = S(n)+S(k)-9r$. 

 Proof :  When you add two numbers, and you do a carry, you are taking away 10 from 1 column, and adding 1 to another column, giving a net loss of 9 to the total.

 Thus, we can see that we need to regroup exactly twice when we add 864. And, the lowest possible n is 299, so let's start from there.

 299 gets three regroups, so we are going to need to take away from digits, and dump the excess in the hundred's place, since the hundreds are going to regroup anyways.

 So, if we take away from the tens digit, we need to take away until we get 2 in the tens digit(since the ones will regroup). So, we get the number 929, which works (929+864 = 1793), but is not the smallest. 

 If we take away from the ones digit, we only have to take away 4, turning the unit's place to 5. 5+4 is 9, so it won't regroup. Dump the ones into the hundred's place, and we get the number $\boxed{695}$

 -AlexLikeMath

 





\textbf{Solution 5}

Although this solution doesn't directly solve the problem, it greatly hastens the bashing process.

 Call the three digits a, b, and c.

 When you add each of 8, 6, and 4 to a, b, and c the resultant will either get smaller or larger, depending on the original number.

 
e.g. If c is 7, then adding 4 will reduce the 7 to a 1, whilst leaving a one for b. 

 If c is 3, then adding 4 will simply add four to the total, and make the 3 a 7.

 
Each of 8, 6, and 4 all can reduce the original number by a certain amount and can increase the original number by a certain amount.

 8 can reduce by 2 for all numbers greater than 1, 6 can reduce numbers by 4, and 4 can reduce numbers by 6.

 
Possibilities:

 -2, +8, (although this becomes obviously impossible later on) -4, +6, -6, +4

 
Also, realize that if the number is reduced, then a one will be carried to the following decimal place on the left, consequently reducing that amount they reduce. It's like a puzzle! Within no time you should find that if you add 4 to c, subtract 4 from b, subtract 1 from a, and leave a 1 in the thousands place, the total is equated to zero. This is optimal because most of the addition is kept to the left, where the effect to real value is less. (e.g. 299 is smaller than 992)

 
Now you have +1, -1, -4, +4 in the decimals, and a VERY fast trial and error gives $\boxed{695}.$

 
-Jackshi2006

 Postscript by Jackshi2006

 When I revisited the problem I realized that you can actually list out every possible number for n. 695 stands out very clearly as the smallest, because the only other possibility is 595, which doesn’t add to 20 to begin with.

 





\textbf{Solution 6 (Fast)}

First, note that to compute $s(m+n)$ (for any positive integers $m$ and $n$), one can simply find the sum of $s(m)$ and $s(n)$ minus 9 times the number of times one regroups when adding $m$ to $n$.  One can see why this is by noticing that if one were to "forget" to regroup, and leave, say, a 10 in  the ones' place, the sum of the digits would be 9 higher than if one did regroup.  Anyway, one can see that the smallest 3-digit number (on AIME, all the answers are integers from 0 to 999) whose digits sum to 20 is 299.  If we add this number to 864, we have to regroup 3 times, so $s(299+864)$ will be $9=9\cdot3-(8+6+4)$ smaller than $s(299)$.  We want this difference to be 0, so we need to find a way to only regroup two times.

 
We now notice that regrouping the hundreds is inevitable, so we must either prevent regrouping the ones or the tens.  Preventing regrouping the tens would require moving many of the tens to the hundreds' place (the ones' place is already full), which is bad when we are trying to minimize the number, but preventing regrouping the ones requires moving fewer ones to the hundreds' place.

 
We find that to preventing regrouping the ones, the ones' place of our number must be at most 5 (a larger number would sum to ten when added to 4).  Because we want to minimize the number of ones we move to the hundreds' place, we leave exactly 5 by moving four ones to the hundreds' place:  $299\to\boxed{695}$.

 





\textbf{Solution 7 (Bad Solution)}

Bashing out modulo $9$ and getting lucky we get that if the $8$ and $6$ carry over when adding $n$ and $864$, that $100a+10b+c \equiv 1+100(a-1)+10(b-4)+c+4 \pmod{9}$ such that $n=100a+10b+c$ and after maximizing $b$ and $c$ such that $c<6$ to not make the $4$ carry over to minimize $a$ we get that $\boxed{695}$ is our answer after confirming there are no lesser solutions.

 Note: If the sum of the digits are equal, they are also the same modulo $9$.~SirAppel

 





\textbf{Solution 8 (not realistic but works)}

Alternatively, you can use python to solve the problem:

 \begin{verbatim}
def sum\_of\_digits(number):  total = 0  while number > 0:    digit = number \% 10    total += digit    number //= 10  return totalfor i in range(1000):    digitsum = sum\_of\_digits(i)    newdigitsum = sum\_of\_digits(i+864)       if digitsum == newdigitsum:        if digitsum == 20:            print(i)
\end{verbatim}





\textbf{Solution 9}

Let $n=100a+10b+c$, then we have $a+b+c=20$. Note that the digits of the number $299+864=1163$ do not sum to $20$. Hence, we make the substitution $a=2+x_{1}$. We discuss the minimal of $n$ in 4 cases based on the digits $b,c$:

 \textbf{Case 1} : $b < 4, c < 6$

 We see that $b+c\le 8$ and  $\max{\{a\}}=9$, contradicting $a+b+c=20$. 

 \textbf{Case 2} : $b \ge 4, c < 6$

 We make the substitution $b=4+x_{2}$. It follows that
 \[n+864=1000+100(x_{1}+1)+10x_{2}+c+4\]Clearly $1+(x_{1}+1)+x_{2}+c+4=20\iff a+b+c=20$. We see that $c\le 5$. To minimize $n$, let $c=5$. We have
 \[x_{1}+x_{2}=9\]We find $x_{2}\le 5$ and $x_{1}\ge 4$, and $a=2+x_{1}, b=4+x_{2}$. Thus, the minimal of $n$ in this case is $695$.

 \textbf{Case 3} : $b < 4, c \ge 6$

 We make the substitution $c=6+x_{3}$. It follows that
 \[n+864=1000+100x_{1}+10(b+7)+x_{3}\]If $b<3$, adding the digits of $n+864$ yields $1+x_{1}+b+7+x_{3}=20\iff a+b+c=20$. We see that $n$ is minimized when $x_{3}=3, b=2$. This gives $x_{1}=7$. The minimal of $n$ in this regard is $929$.

 If $b=3$, then the digits of $n+864$ are $1, x_{1}+1, x_{3}$, which cannot sum to $20$.

 \textbf{Case 4} : $b \ge 4, c \ge 6$

 We make the substitution $b=4+x_{2}, c=6+x_{3}$. It follows that
 \[n+864=1000+100(x_{1}+1)+10(x_{2}+1)+x_{3}\]Yet, by our substitution, $1+(x_{1}+1)+(x_{2}+1)+x_{3}=20\iff a+b+c=29$. This is a contradiction, too.

 Therefore, the minimal of $n$ is $\boxed{695}$.

 




	\item (\textit{AIME I 2015}) Let $S$ be the set of all ordered triple of integers $(a_1,a_2,a_3)$ with $1 \le a_1,a_2,a_3 \le 10$. Each ordered triple in $S$ generates a sequence according to the rule $a_n=a_{n-1}\cdot | a_{n-2}-a_{n-3} |$ for all $n\ge 4$. Find the number of such sequences for which $a_n=0$ for some $n$.


 
 \textbf{Answer: }494



\textbf{Solution 1}

Let $a_1=x, a_2=y, a_3=z$. First note that if any absolute value equals 0, then $a_n=0$.Also note that if at any position, $a_n=a_{n-1}$, then $a_{n+2}=0$.Then, if any absolute value equals 1, then $a_n=0$.Therefore, if either $|y-x|$ or $|z-y|$ is less than or equal to 1, then that ordered triple meets the criteria.Assume that to be the only way the criteria is met.To prove, let $|y-x|>1$, and $|z-y|>1$. Then, $a_4 \ge 2z$, $a_5 \ge 4z$, and $a_6 \ge 4z^2$.However, since the minimum values of $a_5$ and $a_6$ are equal, there must be a scenario where the criteria was met that does not meet our earlier scenarios. Calculation shows that to be $z=1$, $|y-x|=2$. Again assume that any other scenario will not meet criteria.To prove, divide the other scenarios into two cases: $z>1$, $|y-x|>1$, and $|z-y|>1$; and $z=1$, $|y-x|>2$, and $|z-y|>1$.For the first one, $a_4 \ge 2z$, $a_5 \ge 4z$, $a_6 \ge 8z$, and $a_7 \ge 16z$, by which point we see that this function diverges.For the second one, $a_4 \ge 3$, $a_5 \ge 6$, $a_6 \ge 18$, and $a_7 \ge 54$, by which point we see that this function diverges.

 
Therefore, the only scenarios where $a_n=0$ is when any of the following are met:$|y-x|<2$ (280 options)$|z-y|<2$ (280 options, 80 of which coincide with option 1)$z=1$, $|y-x|=2$. (16 options, 2 of which coincide with either option 1 or option 2)Adding the total number of such ordered triples yields $280+280-80+16-2=\boxed{494}$.

 
                \begin{itemize}
                    \item Note to author:
                \end{itemize}


                Because $a_4 \ge 2z$, $a_5 \ge 4z$, $a_6 \ge 8z$, and $a_7 \ge 16z$ doesn't mean the function diverges. What if $z = 7$, $a_4 = 60$, and $a_5 = 30$ too?

 
                \begin{itemize}
                    \item Note to the Note to author:
                \end{itemize}


                This isn't possible because the difference between is either 0, 1, or some number greater than 1. If $a_4 = 63\text{, then }a_5$ is either 0, 63, or some number greater than 63 (since it must be a multiple of $z$).

 





\textbf{Solution 2}

Note that the only way for a $0$ to be produced at $a_n$ is if either $a_{n-1} = 0$ or $a_{n-2} = a_{n-3}$. Since the first one will eventually get to the first three assuming that there is no $a_{n-2} = a_{n-3}$ for any $n$, that is not possible because $a_1 , a_2 , a_3 >= 1$. Therefore, we need $a_{n - 2} = a_{n - 3}$. 

 If $2$ consecutive numbers out of $a_1 , a_2 , a_3$ are equal, then those cases work($a_1$ and $a_2$ or $a_2$ and $a_3$ $\textbf{NOT}$ $a_1$ and $a_3$). This is simply $10 \cdot 10 + 10 \cdot 10 - 10 = 190$ by PIE. 

 Now, note that if any of the first three numbers have difference of $1$, we have another working case. First, we calculate how many there are given exactly one of $a_1,a_2$ or $a_2,a_3$ have difference $1$. Given $3$ numbers such that the first $2$ have difference $1$, exactly $4$ permutations work(assuming the numbers are $x,y,$ and $z$ such that $|x-y| = 1$): $x,y,z$; $y,x,z$; $z,x,y$; and $z,y,x$. If the two consecutive numbers are $1$ and $2$, then the last number has $7$ possiblities: $4,5,6,\cdots , 10$. This is symmetric for $9$ and $10$. If the consecutive numbers are $(2,3),\cdots , (8,9)$, there are $6$ possibilities($10$ minus the numbers themselves and the numbers directly above and below). Note that we are not counting any cases already counted in the first case. Therefore, this case gives you $4(7 + 6 * 7 + 7) = 224$. Now we consider the case that there both adjacent $a$s have increments of $1$. 
 \[+1 , +1 \rightarrow 8\] 
 \[-1 , -1 \rightarrow 8\] 
 \[+1 , -1 \rightarrow 9\] 
 \[-1 , +1 \rightarrow 9\] Therefore this gives $224 + 8 + 8 + 9 + 9 = 258$. However, note that we have to add the case where you have $3$ consecutive numbers in arrangement such that only $2$ consecutive numbers have difference $1$. For example, $1,3,2$ is one such triple. There are $8$ triples of consecutive numbers and $4$ ways to arrange each one(e.x: $(1,3,2);(3,1,2);(2,1,3);(2,3,1)$). This adds on 32 working cases, so this case gives $258 + 32 = 290$.

 Note that there is an ultraspecial(Yes, I know that's not a word) case where we generate a pair of $a_i$ that have difference one. This can only happen if $a_3 = 1$ and $a_1$ and $a_2$ have difference $2$. This contributes $14$ cases($2 * 8$ and then subtract $2$ because of the cases $3,1,1$ and $4,2,1$).

 Therefore, our answer is $190 + 290 + 14 = \boxed{494}$.Solution by hyxue

 





\textbf{Solution 3}

Note that $a_n=0$ for some index $n$ is equivalent to $a_{n-2}=a_{n-3}$. We proceed with casework, where overlap with the former cases are excluded.

 \textbf{Case 1: }$a_{2}=a_{1}$

 There are clearly $10\cdot 10=100$ combinations.

 \textbf{Case 2: }$a_{3}=a_{2}$

 There are clearly $10\cdot 9=90$ combinations.

 \textbf{Case 3: }$a_{4}=a_{3}\iff \lvert a_{2}-a_{1}\lvert =1$

 For $a_{2}=2,3,\ldots, 9$ there are $9$ ways to choose $a_{3}$ and $2$ ways to choose $a_{1}$. For $a_{2}=1, 10$, there are $9$ ways to choose $a_{3}$, and $a_{1}$ is fixed. In total, there are $9\cdot 8\cdot 2+9\cdot 2=162$ combinations.

 \textbf{Case 4: }$a_{5}=a_{4}\iff \lvert a_{3}-a_{2}\lvert =1$

 For $a_{2}=2,3,\ldots, 9$ there are $7$ ways to choose $a_{1}$ and $2$ ways to choose $a_{3}$. For $a_{2}=1, 10$, there are $8$ ways to choose $a_{1}$, and $a_{3}$ is fixed. In total, there are $7\cdot 8\cdot 2+2\cdot 8=128$ combinations.

 \textbf{Case 5: }$a_{6}=a_{5}\iff \lvert a_{4}-a_{3}\lvert =1 \iff a_{3}=1, \lvert a_{2}-a_{1}\rvert = 2$

 Note that for $a_{3}\ge 2$, any multiple $a_{4}=k\cdot a_{3}$ never suffices $\lvert a_{4}-a_{3}\rvert=1$, as the equation $(k-1)a_{3}=1$ has no integer solutions. Hence, we must have $a_{3}=1$ to achieve this case.

 For $a_{2}=3,4,\ldots, 8$, there are $2$ ways to choose $a_{1}$. For $a_{2}=9,10$, $a_{1}$ is fixed. In total, there are $2\cdot 6+2=14$ combinations.

 \textbf{Case 6: }$a_{7}=a_{6}\iff \lvert a_{5}-a_{4}\lvert =1$

 By a similar observation from \textbf{Case 5}, this is equivalent to $a_{4}=1, \lvert a_{3}-a_{2}\rvert=2\iff a_{3}=1, \lvert a_{2}-a_{1}\rvert=1$. This implies our previous conditions is necessary for the cases beyond \textbf{Case 5}. Indeed, proceeding inductively, we see that \begin{array}{c|c|c}\text{Case} & \text{Condition} & \text{Overlap with:} \\ \hline\text{Case 6}  & a_{3}=1, \lvert a_{2}-a_{1}\rvert=1 & \text{Case 3}\\\text{Case 7}  & a_{4}=1, \lvert a_{3}-a_{2}\rvert=1 & \text{Case 4}\\\text{Case 8}  & a_{5}=1, \lvert a_{4}-a_{3}\rvert=1 & \text{Case 5}\\\text{Case 9}  & a_{6}=1, \lvert a_{5}-a_{4}\rvert=1 & \text{Case 3}\\\ldots & \ldots &\ldots\end{array}Thus, all possible sequences are discussed.

 Therefore, the total of such sequences is $100+90+162+128+14=\boxed{494}$, as desired.

 


 




	\item (\textit{AIME I 2015}) Let $f(x)$ be a third-degree polynomial with real coefficients satisfying

 \[|f(1)|=|f(2)|=|f(3)|=|f(5)|=|f(6)|=|f(7)|=12.\] Find $|f(0)|$.


 



 
 \textbf{Answer: }072



\textbf{Solution 1}

Let $f(x)$ = $ax^3+bx^2+cx+d$.Since $f(x)$ is a third degree polynomial, it can have at most two bends in it where it goes from up to down, or from down to up.By drawing a coordinate axis, and two lines representing $12$ and $-12$, it is easy to see that $f(1)=f(5)=f(6)$, and $f(2)=f(3)=f(7)$; otherwise more bends would be required in the graph. Since only the absolute value of $f(0)$ is required, there is no loss of generalization by stating that $f(1)=12$, and $f(2)=-12$. This provides the following system of equations.
 \[a +     b +   c +   d =  12\]
 \[8a +   4b + 2c +   d = -12\]
 \[27a +   9b + 3c +   d = -12\]
 \[125a + 25b + 5c +   d =  12\]
 \[216a + 36b + 6c +   d =  12\]
 \[343a + 49b + 7c +   d = -12\]Using any four of these functions as a system of equations yields $d = |f(0)| = \boxed{072}$

 Note: You can use \href{https://en.wikipedia.org/wiki/Gaussian\_elimination}{Gaussian elimination} to solve these equations.

 





\textbf{Solution 2}

By drawing the function, and similar to Solution 1, WLOG let $f(1)=f(5)=f(6)=12$. Then, $f(2)=f(3)=f(7)$. Set $g(x)+12=f(x)$. Then the roots of $g(x)$ are $1,5,6$. So, $g(x)=a(x-1)(x-5)(x-6)$. Plug in $x=2$ to find a. We know 
 \[-24=-12-12=f(2)-12=g(2)=a(1)(-3)(-4)=12a\].So, $a=-2$. Thus, $f(x)=g(x)+12=-2(x-1)(x-5)(x-6)+12$, and then $|f(0)|=60+12=\boxed{072}$.

 





\textbf{Solution 3}

Without loss of generality, let $f(1) = 12$. (If $f(1) = -12$, then take $-f(x)$ as the polynomial, which leaves $|f(0)|$ unchanged.) Because $f$ is third-degree, write
 \[f(x) - 12 = a(x - 1)(x - b)(x - c)\]
 \[f(x) + 12 = a(x - d)(x - e)(x - f)\]where $\{b, c, d, e, f \}$ clearly must be a permutation of $\{2, 3, 5, 6, 7\}$ from the given condition. Thus $b + c + d + e + f = 2 + 3 + 5 + 6 + 7 = 23.$ However, subtracting the two equations gives $-24 = a[(x - 1)(x - b)(x - c) - (x - d)(x - e)(x - f)]$, so comparing $x^2$ coefficients gives $1 + b + c = d + e + f$ and thus both values equal to $\dfrac{24}{2} = 12$. As a result, $\{b, c \} = \{5, 6 \}$. As a result, $-24 = a (12)$ and so $a = -2$. Now, we easily deduce that $f(0) = (-2) \cdot (-1) \cdot (-5) \cdot (-6) + 12 = 72,$ and so removing the without loss of generality gives $|f(0)| = \boxed{072}$, which is our answer.

 


 





\textbf{Solution 4}

The following solution is similar to solution 3, but assumes nothing. Let $g(x)=(f(x))^2-144$. Since $f$ has degree 3, $g$ has degree 6 and has roots 1,2,3,5,6, and 7. Therefore, $g(x)=k(x-1)(x-2)(x-3)(x-5)(x-6)(x-7)$ for some $k$. Hence $|f(0)|=\sqrt{g(0)+144}=\sqrt{1260k+144}$. Note that $g(x)=(f(x)+12)(f(x)-12)$. Since $f$ has degree 3, so do $f(x)+12$ and $f(x)-12$; and both have the same leading coefficient. Hence $f(x)+12=a(x-q)(x-r)(x-s)$ and $f(x)-12=a(x-t)(x-u)(x-v)$ for some $a\neq 0$ (else $f$ is not cubic) where $\{q,r,s,t,u,v\}$ is the same as the set $\{1,2,3,5,6,7\}$. Subtracting the second equation from the first, expanding, and collecting like terms, we have that 
 \[24=a((t+u+v-(q+r+s))x^2-a(tu+uv+tv-(qr+qs+rs))x+a(tuv-qrs)\]which must hold for all $x$. Since $a\neq 0$ we have that (1) $t+u+v=q+r+s$, (2) $tu+uv+tv=qr+qs+rs$ and (3) $a(tuv-qrs)=24$. Since $q+r+s+t+u+v$ is the sum of 1,2,3,5,6, and 7, we have $q+r+s+t+u+v=24$ so that by (1) we have $q+r+s=12$ and $t+u+v=12$. We must partition 1,2,3,5,6,7 into 2 sets each with a sum of 12. Consider the set that contains 7. It can't contain 6 or 5 because the sum of that set would already be $\geq 12$ with only 2 elements. If 1 is in that set, the other element must be 4 which is impossible. Hence the two sets must be $\{2,3,7\}$ and $\{1,5,6\}$. Note that each of these sets happily satisfy (2). By (3), since the sets have products 42 and 30 we have that $|a|=\frac{24}{|tuv-qrs|}=\frac{24}{12}=2$. Since $a$ is the leading coefficient of $f(x)$, the leading coefficient of $(f(x))^2$ is $a^2=|a|^2=2^2=4$. Thus the leading coefficient of $g(x)$ is 4, i.e. $k=4$. Then from earlier, $|f(0)|=\sqrt{g(0)+144}=\sqrt{1260k+144}=\sqrt{1260\cdot4+144}=\sqrt{5184}=72$ so that the answer is $\boxed{072}$.

 


 





\textbf{Solution 5}

Express $f(x)$ in terms of powers of $(x-4)$:
 \[f(x) = a(x-4)^3 + b(x-4)^2 + c(x-4) + d\]By the same argument as in the first Solution, we see that $f(x)$ is an odd function about the line $x=4$, so its coefficients $b$ and $d$ are 0. From there it is relatively simple to solve $f(2)=f(3)=-12$ (as in the above solution, but with a smaller system of equations):
 \[a(1)^3 + c(1) = -12\]
 \[a(2)^3 + c(2) = -12\]$a=2$ and $c=-14$
 \[|f(0)| = |2(-4)^3 - 14(-4)| = \boxed{072}\]

 


 





\textbf{Solution 6 (Finite differences)}

Because a cubic must come in a "wave form" with two points of inflection, we can see that $f(1)=f(5)=f(6)$, and $f(2)=f(3)=f(7)$. By symmetry, $f(4)=0$. Now, WLOG let $f(1)=12$, and $f(2)=f(3)=-12$. Then, we can use finite differences to get that the third (constant) difference is $12$, and therefore $f(0)=12+(24+(24+12))=\boxed{072}$.

 





\textbf{Solution 7 (Like solution 1 without annoying systems)}

We can rewrite our function as two different cubics, $f(x)=k(x-a)(x-b)(x-c)\pm12=k(x-d)(x-e)(x-f)\mp12$. Note that $k$ is the same in both because they must be equal, so their leading terms must be. We must then, following Vieta's, choose our roots such that $a+b+c=d+e+f$ and verify that the other Vieta's formulas hold. Additionally, a cubic must only cross the x-axis thrice, restricting our choices for roots further. Choosing $a=1$, $b=5$, $c=6$, $d=2$, $e=3$, $f=7$ yields: 
 \[kx^3-12kx^2+41kx-30k\pm12=kx^3-12kx^2+41kx-42k\mp12\] For the constant terms to have a difference of 24 ($|\pm12-\mp12|$), we must have $k=\pm2$, so the constant term of our polynomial is $\pm72$, the absolute value of which is $\boxed{072}$. -- Solution by eiis1000

 


 





\textbf{Solution 8 (First few steps of solution 1)}

We can rewrite the function as $f(x) - 12 = z(x-1)(x-5)(x-6)$ and $f(x) + 12 = z(x-2)(x-3)(x-7)$. Since we need to find $|f(0)|$, substitute 0 for x in these two equations. 

 
 \[f(0) - 12 = z(0-1)(0-5)(0-6), f(0) + 12 = z(0-2)(0-3)(0-7)\] 
 \[f(0) - 12 = z\cdot -30, f(0) + 12 = z\cdot -42\]Isolating $z$ in both of the equations, 
 \[z = \frac{f(0)-12}{-30}, z = \frac{f(0) + 12}{-42}\]Equating the two and solving for $f(0)$, we see $|f(0)| = 72$.

 





\textbf{Solution 9 (Cheese Solution)}

Let the leading coefficient of  $f(x)$ be $a$. Then, it is obvious that $(f(x) - 12)(f(x) + 12) = a^2(x - 1)(x - 2)(x - 3)(x - 5)(x - 6)(x - 7)$. Let us now, let $x = 0$. We then have, after cleaning it up nicely, $[f(0)]^2 = 2^2 * 3^2 * (35a^2 + 4)$. We now take the square root of both sides, to obtain $|f(0)|$. So, $|f(0)| = 6 * \sqrt{35a^2 + 4}$. Now, this is the cheese part. Since this is aime, we know that the answer must be an integer, so we assume $a = 2$. Thus, we get $|f(0)| = 6 * \sqrt{144} = 6 * 12 = 072$.

 





\textbf{Solution 10 (Rigor Olympiad Analysis)}

Without loss of generality, assume $f(x)$ is monic.

 \textbf{Lemma 1: For global minimums $a$, $b$ we have $a < b$ and increase of the function $f(x)$ upon $(-\infty, a) \cup (b, \infty)$}

 \textbf{Proof of the Lemma 1:}

 Definition of cubic functions show $a$ and $b$ exist. Increase due to monoicy. 

 \textbf{Lemma 2: For global minimums $a,b$ we have $a<b$ and decrease of the function $f(x)$ upon $(a,b)$}.

 \textbf{Proof of Lemma 2:}

 \textbf{Lemma 1} proves the very existence of $a,b$. Definition of global minimums within cubic polynomials directly prove the lemma.

 \textbf{Lemma 3: $f(1) = f(5) = f(6) =- 12$ and the remaining values refer 12.}

 \textbf{Proof of Lemma 3:}

 Firstly, cubic functions are always continuous. If we wish to arrange such a configuration we have to place two points $a$ and $b$, with $a < b$ on the plane such that the function intersects the lines $y = 12, -12$ at 6 specific points. A cubic intersects two horizontal lines at 6 points. Surprisingly, there are 6 points given to us. At maximum a horizontal line is intersected twice, so the pattern: $f(1) = -12, f(2) = 12, f(3) = 12, f(5) = -12, f(6) = -12, f(7) = 12$ is obtained, proving the lemma.

 \textbf{2015 AIME I Q10}

 Let $f(x) = c(x-1)(x-5)(x-6) + 12$ and $f(x) = c(x-2)(x-3)(x-7) - 12$. Then, we have $30c + 12 = 42c - 12$, so $c = 2$. Then, $f(0) = 2(30) + 12 = 60 + 12 = \boxed{72}$.

 




	\item (\textit{AIME I 2015}) Triangle $ABC$ has positive integer side lengths with $AB=AC$. Let $I$ be the intersection of the bisectors of $\angle B$ and $\angle C$. Suppose $BI=8$. Find the smallest possible perimeter of $\triangle ABC$.


 



 
 \textbf{Answer: }108



\textbf{Solution 1}

Let $D$ be the midpoint of $\overline{BC}$. Then by SAS Congruence, $\triangle ABD \cong \triangle ACD$, so $\angle ADB = \angle ADC = 90^o$.

 Now let $BD=y$, $AB=x$, and $\angle IBD = \dfrac{\angle ABD}{2} = \theta$.

 Then $\mathrm{cos}{(\theta)} = \dfrac{y}{8}$

 and $\mathrm{cos}{(2\theta)} = \dfrac{y}{x} = 2\mathrm{cos^2}{(\theta)} - 1 = \dfrac{y^2-32}{32}$.

 Cross-multiplying yields $32y = x(y^2-32)$.

 Since $x,y>0$, $y^2-32$ must be positive, so $y > 5.5$.

 Additionally, since $\triangle IBD$ has hypotenuse $\overline{IB}$ of length $8$, $BD=y < 8$.

 Therefore, given that $BC=2y$ is an integer, the only possible values for $y$ are $6$, $6.5$, $7$, and $7.5$.

 However, only one of these values, $y=6$, yields an integral value for $AB=x$, so we conclude that $y=6$ and $x=\dfrac{32(6)}{(6)^2-32}=48$.

 Thus the perimeter of $\triangle ABC$ must be $2(x+y) = \boxed{108}$.

 





\textbf{Solution 2 (No Trig)}

Let $AB=x$ and the foot of the altitude from $A$ to $BC$ be point $E$ and $BE=y$. Since ABC is isosceles, $I$ is on $AE$. By Pythagorean Theorem, $AE=\sqrt{x^2-y^2}$. Let $IE=a$ and $IA=b$. By Angle Bisector theorem, $\frac{y}{a}=\frac{x}{b}$. Also, $a+b=\sqrt{x^2-y^2}$. Solving for $a$, we get $a=\frac{\sqrt{x^2-y^2}}{1+\frac{x}{y}}$. Then, using Pythagorean Theorem on $\triangle BEI$ we have $y^2+\left(\frac{\sqrt{x^2-y^2}}{1+\frac{x}{y}}\right)^2=8^2=64$. Simplifying, we have $y^2+y^2\frac{x^2-y^2}{(x+y)^2}=64$. Factoring out the $y^2$, we have $y^2\left(1+\frac{x^2-y^2}{(x+y)^2}\right)=64$. Adding 1 to the fraction and simplifying, we have $\frac{y^2x(x+y)}{(x+y)^2}=32$. Crossing out the $x+y$, and solving for $x$ yields $32y = x(y^2-32)$. Then, we continue as Solution 1 does.

 





\textbf{Solution 3}

Let $AB=x$, call the midpoint of $BC$ point $E$, call the point where the incircle meets $AB$ point $D$, 

 and let $BE=y$. We are looking for the minimum value of $2(x+y)$. $AE$ is an altitude because the triangle

 is isosceles. By Pythagoras on $BEI$, the inradius is $\sqrt{64-y^2}$ and by Pythagoras on $ABE$, $AE$ is

 $\sqrt{x^2-y^2}$. By equal tangents, $BE=BD=y$, so $AD=x-y$. Since $ID$ is an inradius, $ID=IE$ andusing pythagoras on $ADI$ yields $AI=$$\sqrt{x^2-2xy+64}$. $ADI$ is similar to $AEB$ by $AA$, so we 

 can write $\frac{x-y}{\sqrt{x^2-2xy+64}}=\frac{\sqrt{x^2-y^2}}{x}$. Simplifying, $\frac{x}{\sqrt{x^2-2xy+64}}=\sqrt{\frac{x+y}{x-y}}$. 

 Squaring, subtracting 1 from both sides, and multiplying everything out, we get $yx^2-2xy^2+64y=yx^2 -32x+32y-xy^2$, which turns into $32y=x(y^2-32)$. Finish as in Solution 1.

 





\textbf{Solution 4}

Angle bisectors motivate trig bash. Define angle $IBC = x$. Foot of perpendicular from $I$ to $BC$ is point $P$. $\overline{BC} = 2\overline{BP} = 2(8\cos(x)) = N$, where $N$ is an integer. Thus, $\cos(x) = \frac{N}{16}$. Via double angle, we calculate $\overline{AB}$ to be $\frac{8\cos(x)}{2\cos(x)^2 - 1} = \frac{64N}{N^2 - 128}$. This is to be an integer.  We can bound $N$ now, as $N > 11$ to avoid negative values and $N < 16$ due to triangle inequality. Testing, $N = 12$ works, giving $\overline{AB} = 48, \overline{BC} = 12$.Our answer is $2 * 48 + 12 = \boxed{108}$.- whatRthose

 





\textbf{Solution 5}


                    \begin{figure}[H]
                        \centering
                        \includegraphics[width=0.5\textwidth]{Solution/AIME_2015_I_11_Solution_Figure_1.png}
                    \end{figure}
                    Let $M$ be midpoint $BC, BM = x, AB = y, \angle IBM = \alpha.$

 $BI$ is the bisector of $\angle ABM$ in $\triangle ABM.$$BI = \frac {2 xy \cos \alpha}{x+y} = 8, \cos \alpha = \frac {x}{8} \implies \frac {x^2 y}{x+y} = 32.$
 \[y = \frac {32 x} {x^2 - 32}.\]$BC = 2x$ is integer, $5.5^2 < 32 \implies x \ge 6.$$BM < BI \implies x =\{ 6, 6.5, 7, 7.5 \}.$

 If $x > 6$ then $y$ is not integer. 
 \[x = 6 \implies y = 48 \implies 2(x+y) =  \boxed{\textbf{108}}.\]\textbf{vladimir.shelomovskii@gmail.com, vvsss}

 


 


 


 




	\item (\textit{AIME I 2015}) Consider all 1000-element subsets of the set $\{ 1, 2, 3, ... , 2015 \}$.  From each such subset choose the least element.  The arithmetic mean of all of these least elements is $\frac{p}{q}$, where $p$ and $q$ are relatively prime positive integers.  Find $p + q$.


 



 
 \textbf{Answer: }431



\textbf{Solution 1}

Let $M$ be the desired mean. Then because $\dbinom{2015}{1000}$ subsets have 1000 elements and $\dbinom{2015 - i}{999}$ have $i$ as their least element,
 \begin{align*} \binom{2015}{1000} M &= 1 \cdot \binom{2014}{999} + 2 \cdot \binom{2013}{999} + \dots + 1016 \cdot \binom{999}{999} \\      &= \binom{2014}{999} + \binom{2013}{999} + \dots + \binom{999}{999} \\      &                      + \binom{2013}{999} + \binom{2012}{999} + \dots + \binom{999}{999} \\      & \dots \\      & + \binom{999}{999} \\      &= \binom{2015}{1000} + \binom{2014}{1000} + \dots + \binom{1000}{1000} \\      &= \binom{2016}{1001}. \end{align*}
Using the definition of binomial coefficient and the identity $n! = n \cdot (n-1)!$, we deduce that
 \[M = \frac{2016}{1001} = \frac{288}{143}.\]The answer is $\boxed{431}.$

 \textbf{Potential Inspiration} For solution 1, the inspiration could be that once we select the set $\{a_1, a_2,...,a_{2015}\}$, where $a_1<a_2<...<a_{2015}$, we want to multiply each such set by $a_1$ and sum up through all such sets.

 So, how do we scale each such set by $a_1$? We realize that if we were to choose one more element, specifically, less than $a_1$, this could be done in $a_1-1$ ways.

 Oh! So, now if we add $1$ to all the numbers in our set, they become $\{a_1+1, a_2+1,...,a_{2015}+1\}$, so the number of ways to pick another number below the least number in this set is $a_1$!

 





\textbf{Solution 2}

Each 1000-element subset $\left\{ a_1, a_2,a_3,...,a_{1000}\right\}$ of $\left\{1,2,3,...,2015\right\}$ with $a_1<a_2<a_3<...<a_{1000}$ contributes $a_1$ to the sum of the least element of each subset. Now, consider the set $\left\{a_1+1,a_2+1,a_3+1,...,a_{1000}+1\right\}$. There are $a_1$ ways to choose a positive integer $k$ such that $k<a_1+1<a_2+1,a_3+1<...<a_{1000}+1$ ($k$ can be anything from $1$ to $a_1$ inclusive). Thus, the number of ways to choose the set $\left\{k,a_1+1,a_2+1,a_3+1,...,a_{1000}+1\right\}$ is equal to the sum. But choosing a set $\left\{k,a_1+1,a_2+1,a_3+1,...,a_{1000}+1\right\}$ is the same as choosing a 1001-element subset from $\left\{1,2,3,...,2016\right\}$! 

 Thus, the average is $\frac{\binom{2016}{1001}}{\binom{2015}{1000}}=\frac{2016}{1001}=\frac{288}{143}$. Our answer is $p+q=288+143=\boxed{431}$.

 





\textbf{Solution 3(NOT RIGOROUS)}

Let $p$ be the size of the large set and $q$ be the size of the subset (i.e. in this problem, $p = 2015$ and $q = 1000$). We can easily find the answers for smaller values of $p$ and $q$:

 For $p = 2$ and $q = 2$, the answer is $1$.

 For $p = 3$ and $q = 2$, the answer is $\frac43$.

 For $p = 4$ and $q = 2$, the answer is $\frac53$.

 For $p = 3$ and $q = 3$, the answer is $1$.

 For $p = 4$ and $q = 3$, the answer is $\frac54$.

 For $p = 5$ and $q = 3$, the answer is $\frac32$.

 At this point, we can see a pattern: our desired answer is always $\frac{p+1}{q+1}$. Plugging in $p = 2015$ and $q = 1000$, the answer is $\frac{2016}{1001}=\frac{288}{143}$, so $288 + 143 = \boxed{431}$.

 





\textbf{Solution 4 (short)}

In the "average case", the numbers evenly partition the interval [0,2016] into 1001 parts. Then because it asks for the expected value of the least element the answer is $\frac{2016}{1001}$.

 -tigershark22

 VIEWER NOTE: This solution doesn't always work, for example, take $n^2$ on $[0,1]$. The "average case" is $n=\frac{1}{2}\implies n^2=\frac{1}{4}$ but integrating $n^2$ from $0$ to $1$ we see that definitely is not the case. This works only in certain situations, so maybe this solution would be better off with proof that this is one of those situations.

 





\textbf{Solution 5}

When it says average, it really just means expected value. Let's do casework.

 Our first case is when the smallest number in the subset is $1$. There are $2014$ numbers to choose from, and $999$ numbers left to choose. So, this contributes $\dbinom{2014}{999}$ to our expected value.

 Then, for the smallest number being $2$, we have $2013$ numbers to choose from and $999$ numbers left to choose. So, this contributes $2\dbinom{2013}{999}$.

 Similarly, the smallest number being $3$ contributes $3\dbinom{2012}{999}$ to our expected value. We keep this up all the way to $1016\dbinom{999}{999}$.

 This makes our sum
 \[\dbinom{2014}{999} + 2\dbinom{2013}{999} + 3\dbinom{2012}{999} + \cdots + 1016\dbinom{999}{999}.\]

 To evaluate this, we can split each binomial into groups of $1$, like this:
 \[\binom{2014}{999} + \binom{2013}{999} + \binom{2013}{999} + \binom{2012}{999} + \binom{2012}{999} + \binom{2012}{999} + \cdots + \underbrace{\binom{999}{999} + \binom{999}{999} + \cdots + \binom{999}{999}}_{\text{1016 }\binom{999}{999}\text{'s}}.\]

 Now, we can turn this into:
 \[\left( \binom{2014}{999} + \binom{2013}{999} + \cdots + \binom{999}{999} \right) + \left( \binom{2013}{999} + \binom{2012}{999} + \cdots + \binom{999}{999} \right) + \cdots + \left( \binom{1000}{999} + \binom{999}{999} \right) + \binom{999}{999}.\]

 Now, by the Hockey Stick Identity, we can further simplify this into
 \[\dbinom{2015}{1000} + \dbinom{2014}{1000} + \cdots + \dbinom{1001}{1000} + \dbinom{1000}{1000}.\]

 Further simplifying this, we get
 \[\dbinom{2016}{1001}.\]

 If this is our sum and there are $\dbinom{2015}{1000}$ possible subsets, our average is simply
 \[\dfrac{\dbinom{2016}{1001}}{\dbinom{2015}{1000}} = \dfrac{2016!}{1001!\,1015!} \times \dfrac{1000!\,1015!}{2015!} = \dfrac{2016}{1001} = \dfrac{288}{143} \Longrightarrow \boxed{431}.\]

 -jb2015007

 \textbf{Generalization} \href{https://artofproblemsolving.com/wiki/index.php/1981\_IMO\_Problems/Problem\_2}{https://artofproblemsolving.com/wiki/index.php/1981\_IMO\_Problems/Problem\_2}

 




	\item (\textit{AIME I 2015}) With all angles measured in degrees, the product $\prod_{k=1}^{45} \csc^2(2k-1)^\circ=m^n$, where $m$ and $n$ are integers greater than 1. Find $m+n$.


 



 
 \textbf{Answer: }091



\textbf{Solution 1}

Let $x = \cos 1^\circ + i \sin 1^\circ$. Then from the identity
 \[\sin 1 = \frac{x - \frac{1}{x}}{2i} = \frac{x^2 - 1}{2 i x},\]we deduce that (taking absolute values and noticing $|x| = 1$)
 \[|2\sin 1| = |x^2 - 1|.\]But because $\csc$ is the reciprocal of $\sin$ and because $\sin z = \sin (180^\circ - z)$, if we let our product be $M$ then
 \[\frac{1}{M} = \sin 1^\circ \sin 3^\circ \sin 5^\circ \dots \sin 177^\circ \sin 179^\circ\]
 \[= \frac{1}{2^{90}} |x^2 - 1| |x^6 - 1| |x^{10} - 1| \dots |x^{354} - 1| |x^{358} - 1|\]because $\sin$ is positive in the first and second quadrants. Now, notice that $x^2, x^6, x^{10}, \dots, x^{358}$ are the roots of $z^{90} + 1 = 0.$ Hence, we can write $(z - x^2)(z - x^6)\dots (z - x^{358}) = z^{90} + 1$, and so
 \[\frac{1}{M} = \dfrac{1}{2^{90}}|1 - x^2| |1 - x^6| \dots |1 - x^{358}| = \dfrac{1}{2^{90}} |1^{90} + 1| = \dfrac{1}{2^{89}}.\]It is easy to see that $M = 2^{89}$ and that our answer is $2 + 89 = \boxed{091}$.

 





\textbf{Solution 2}

Let $p=\sin1\sin3\sin5...\sin89$

 
 \[p=\sqrt{\sin1\sin3\sin5...\sin177\sin179}\]

 
 \[=\sqrt{\frac{\sin1\sin2\sin3\sin4...\sin177\sin178\sin179}{\sin2\sin4\sin6\sin8...\sin176\sin178}}\]

 
 \[=\sqrt{\frac{\sin1\sin2\sin3\sin4...\sin177\sin178\sin179}{(2\sin1\cos1)\cdot(2\sin2\cos2)\cdot(2\sin3\cos3)\cdot....\cdot(2\sin89\cos89)}}\]

 
 \[=\sqrt{\frac{1}{2^{89}}\frac{\sin90\sin91\sin92\sin93...\sin177\sin178\sin179}{\cos1\cos2\cos3\cos4...\cos89}}\]

 $=\sqrt{\frac{1}{2^{89}}}$ because of the identity $\sin(90+x)=\cos(x)$

 we want $\frac{1}{p^2}=2^{89}$

 Thus the answer is $2+89=\boxed{091}$

 





\textbf{Solution 3}

Similar to Solution $2$, so we use $\sin{2\theta}=2\sin\theta\cos\theta$ and we find that:
 \begin{align*}\sin(4)\sin(8)\sin(12)\sin(16)\cdots\sin(84)\sin(88)&=(2\sin(2)\cos(2))(2\sin(4)\cos(4))(2\sin(6)\cos(6))(2\sin(8)\cos(8))\cdots(2\sin(42)\cos(42))(2\sin(44)\cos(44))\\ &=(2\sin(2)\sin(88))(2\sin(4))\sin(86))(2\sin(6)\sin(84))(2\sin(8)\sin(82))\cdots(2\sin(42)\sin(48))(2\sin(44)\sin(46))\\ &=2^{22}(\sin(2)\sin(88)\sin(4)\sin(86)\sin(6)\sin(84)\sin(8)\sin(82)\cdots\sin(42)\sin(48)\sin(44)\sin(46))\\ &=2^{22}(\sin(2)\sin(4)\sin(6)\sin(8)\cdots\sin(82)\sin(84)\sin(86)\sin(88))\end{align*}
Now we can cancel the sines of the multiples of $4$:
 \[1=2^{22}(\sin(2)\sin(6)\sin(10)\sin(14)\cdots\sin(82)\sin(86))\]So $\sin(2)\sin(6)\sin(10)\sin(14)\cdots\sin(82)\sin(86)=2^{-22}$ and we can apply the double-angle formula again:
 \begin{align*}2^{-22}&=(\sin(2)\sin(6)\sin(10)\sin(14)\cdots\sin(82)\sin(86)\\ &=(2\sin(1)\cos(1))(2\sin(3)\cos(3))(2\sin(5)\cos(5))(2\sin(7)\cos(7))\cdots(2\sin(41)\cos(41))(2\sin(43)\cos(43))\\ &=(2\sin(1)\sin(89))(2\sin(3)\sin(87))(2\sin(5)\sin(85))(2\sin(7)\sin(87))\cdots(2\sin(41)\sin(49))(2\sin(43)\sin(47))\\ &=2^{22}(\sin(1)\sin(89)\sin(3)\sin(87)\sin(5)\sin(85)\sin(7)\sin(83)\cdots\sin(41)\sin(49)\sin(43)\sin(47))\\ &=2^{22}(\sin(1)\sin(3)\sin(5)\sin(7)\cdots\sin(41)\sin(43))(\sin(47)\sin(49)\cdots\sin(83)\sin(85)\sin(87)\sin(89))\end{align*}
Of course, $\sin(45)=2^{-\frac{1}{2}}$ is missing, so we multiply it to both sides:
 \[2^{-22}\sin(45)=2^{22}(\sin(1)\sin(3)\sin(5)\sin(7)\cdots\sin(41)\sin(43))(\sin(45))(\sin(47)\sin(49)\cdots\sin(83)\sin(85)\sin(87)\sin(89))\]
 \[\left(2^{-22}\right)\left(2^{-\frac{1}{2}}\right)=2^{22}(\sin(1)\sin(3)\sin(5)\sin(7)\cdots\sin(83)\sin(85)\sin(87)\sin(89))\]
 \[2^{-\frac{45}{2}}=2^{22}(\sin(1)\sin(3)\sin(5)\sin(7)\cdots\sin(83)\sin(85)\sin(87)\sin(89))\]Now isolate the product of the sines:
 \[\sin(1)\sin(3)\sin(5)\sin(7)\cdots\sin(83)\sin(85)\sin(87)\sin(89)=2^{-\frac{89}{2}}\]And the product of the squares of the cosecants as asked for by the problem is the square of the inverse of this number:
 \[\csc^2(1)\csc^2(3)\csc^2(5)\csc^2(7)\cdots\csc^2(83)\csc^2(85)\csc^2(87)\csc^2(89)=\left(\frac{1}{2^{-\frac{89}{2}}}\right)^2=\left(2^{\frac{89}{2}}\right)^2=2^{89}\]The answer is therefore $m+n=(2)+(89)=\boxed{091}$.

 





\textbf{Solution 4}

Let $p=\prod_{k=1}^{45} \csc^2(2k-1)^\circ$.

 Then, $\sqrt{\frac{1}{p}}=\prod_{k=1}^{45} \sin(2k-1)^\circ$.

 Since $\sin\theta=\cos(90^{\circ}-\theta)$, we can multiply both sides by $\frac{\sqrt{2}}{2}$ to get $\sqrt{\frac{1}{2p}}=\prod_{k=1}^{23} \sin(2k-1)^\circ\cos(2k-1)^\circ$.

 Using the double-angle identity $\sin2\theta=2\sin\theta\cos\theta$, we get $\sqrt{\frac{1}{2p}}=\frac{1}{2^{23}}\prod_{k=1}^{23} \sin(4k-2)^\circ$.

 Note that the right-hand side is equal to $\frac{1}{2^{23}}\prod_{k=1}^{45} \sin(2k)^\circ\div \prod_{k=1}^{22} \sin(4k)^\circ$, which is equal to $\frac{1}{2^{23}}\prod_{k=1}^{45} \sin(2k)^\circ\div \prod_{k=1}^{22} 2\sin(2k)^\circ\cos(2k)^\circ$, again, from using our double-angle identity.

 Putting this back into our equation and simplifying gives us $\sqrt{\frac{1}{2p}}=\frac{1}{2^{45}}\prod_{k=23}^{45} \sin(2k)^\circ\div \prod_{k=1}^{22} \cos(2k)^\circ$.

 Using the fact that $\sin\theta=\cos(90^{\circ}-\theta)$ again, our equation simplifies to $\sqrt{\frac{1}{2p}}=\frac{\sin90^\circ}{2^{45}}$, and since $\sin90^\circ=1$, it follows that $2p = 2^{90}$, which implies $p=2^{89}$. Thus, $m+n=2+89=\boxed{091}$.

 





\textbf{Solution 5}

Once we have the tools of complex polynomials there is no need to use the tactical tricks. Everything is so basic (I think).

 Recall that the roots of $x^n+1$ are $e^{\frac{(2k-1)\pi i}{n}}, k=1,2,...,n$, we have
 \[x^n + 1 = \prod_{k=1}^{n}(x-e^{\frac{(2k-1)\pi i}{n}})\]Let $x=1$, and take absolute value of both sides,  
 \[2 = \prod_{k=1}^{n}|1-e^{\frac{(2k-1)\pi i}{n}}|= 2^n\prod_{k=1}^{n}|\sin\frac{(2k-1)\pi}{2n}|\]or, 
 \[\prod_{k=1}^{n}|\sin\frac{(2k-1)\pi}{2n}| = 2^{-(n-1)}\]Let $n$ be even, then,
 \[\sin\frac{(2k-1)\pi}{2n} = \sin\left(\pi - \frac{(2k-1)\pi}{2n}\right) = \sin\left(\frac{(2(n-k+1)-1)\pi}{2n}\right)\]so,
 \[\prod_{k=1}^{n}\left|\sin\frac{(2k-1)\pi}{n}\right| = \prod_{k=1}^{\frac{n}{2}}\sin^2\frac{(2k-1)\pi}{2n}\]Set $n=90$ and we have
 \[\prod_{k=1}^{45}\sin^2\frac{(2k-1)\pi}{180} = 2^{-89}\],
 \[\prod_{k=1}^{45}\csc^2\frac{(2k-1)\pi}{180} = 2^{89}\]-Mathdummy

 





\textbf{Solution 6}

Recall that $\sin\alpha\cdot \sin(60^{\circ}-\alpha)\cdot \sin(60^{\circ}+\alpha)=\frac{1}{4}\cdot \sin3\alpha$Since it is in csc, we can write in sin and then take reciprocal.We can group them by threes, $P=(\sin1^{\circ}\cdot \sin59^{\circ}\cdot \sin61^{\circ})\cdots(\sin29^{\circ}\cdot \sin31^{\circ}\cdot \sin89^{\circ})$. Thus
 \begin{align*} P &=\frac{1}{4^{15}}\cdot \sin3^{\circ}\cdot \sin9^{\circ}\cdots\sin87^{\circ}\\ &=\frac{1}{4^{20}}\cdot \sin9^{\circ}\cdot \sin27^{\circ}\cdot \sin45^{\circ}\cdot \sin63^{\circ}\cdot \sin81^{\circ}\\ &=\frac{1}{4^{20}}\cdot \frac{\sqrt{2}}{2}\cdot \sin9^{\circ}\cdot \cos9^{\circ}\cdot \sin27^{\circ}\cdot \cos27^{\circ}\\ &=\frac{1}{4^{21}}\cdot \frac{\sqrt{2}}{2}\cdot \sin18^{\circ}\cdot \cos36^{\circ}=\frac{\sqrt{2}}{2^{45}} \end{align*}
So we take reciprocal, $\frac 1P=2^{\frac{89}{2}}$, the desired answer is $\frac{1}{P^2}=2^{89}$ leads to answer $\boxed{091}$

 





\textbf{Solution 7}

We have

 
 \[\prod_{k=1}^{45} \csc^2(2k-1)^\circ = \left(\frac{1}{\sin1^\circ \cdot \sin3^\circ \cdots \sin89^\circ}\right)^2.\]

 Multiplying by $\frac{\sin2^\circ \cdot \sin4^\circ \cdots \sin88^\circ}{\sin2^\circ \cdot \sin4^\circ \cdots \sin88^\circ}$ gives

 
 \[\left(\frac{\sin2^\circ \cdot \sin4^\circ \cdots \sin88^\circ}{\sin1^\circ \sin2^\circ \cdot \sin3^\circ \cdots \sin88^\circ \cdot \sin89^\circ}\right)^2\]

 
 \[= \left(\frac{\sin2^\circ \cdot \sin4^\circ \cdots \sin88^\circ}{\sin1^\circ \sin2^\circ \cdot \sin3^\circ \cdots \sin 45^\circ \cdot \cos 44^\circ \cdot \cos 43^\circ \cdots \cos1^\circ}\right)^2.\]

 Using $\sin\alpha \cos\alpha = \frac{1}{2}\sin{2\alpha}$ gives

 
 \[\left(\frac{\sin2^\circ \cdot \sin4^\circ \cdots \sin88^\circ}{\frac{1}{2} \sin2^\circ \cdot \frac{1}{2} \sin4^\circ \cdots \frac{1}{2} \sin88^\circ \cdot \sin45^\circ}\right) ^2\]

 
 \[= \left(\frac{1}{(\frac{1}{2})^{44} \cdot \frac{\sqrt{2}}{2}}\right)^2\]

 
 \[= 2^{89}.\]

 Thus, the answer is $2+89 = \boxed{091}.$

 





\textbf{Solution 8}

Consider the product $\prod_{k=1}^{45} \csc^2(2k-1) = \prod_{k=45}^{1} \sec^2(2k-1) = \prod_{k=1}^{45} \sec^2(2k-1) =  \prod_{k=1}^{45} (1+\tan^2(2k-1))$ However, note that the $45$ numbers in the form $\sqrt{\tan^2(2k-1)} = \tan(2k-1)$ for $1\le{k}\le{45}$ are precisely the roots of the equation $\frac{1}{\tan{(\tan^{-1}{90x}})} = 0$. Thus, it suffices to find $|P(-1)|$, where $P$ is the polynomial formed by the denominator of $\tan{(\tan^{-1}{90\sqrt{x}})}$. This is true because $\prod_{k=1}^{45} (x-\tan^2(2k-1)) = \prod_{k=1}^{45} -(-x+\tan^2(2k-1))$ gives us the factored root form of the equation where $P$ is undefined, which corresponds to the equation where the denominator equals $0$. 

 It remains to find the denominator of $P$; fortunately, we may use tangent angle multiplication rules. Specifically, the denominator of $P$ will be $\sum_{n=0} ^{45} (-1)^n\binom{90}{2n}\sqrt{x}^{2n} = \sum_{n=0} ^{45} (-1)^n\binom{90}{2n}x^n$. Evaluating at $x = -1$, we may obtain the sum $\sum_{n=0} ^{45} (-1)^n\binom{90}{2n}(-1)^{n} = \sum_{n=0} ^{45} \binom{90}{2n}$. 

 From here, there are two ways to finish; the first is to recognize the well known sum $\sum_{n=0} ^{k} \binom{2k}{2n} = 2^(2k-1)$, from which we may plug in $k = 45$ to see $|P(-1)| = 2^(45*2-1) = 2^{89}$ to obtain the answer of $2+89=\boxed{091}$. Otherwise, you may split the previous sum; $\sum_{n=0} ^{45} \binom{90}{2n} = \sum_{n=0} ^{45} \binom{89}{2n-1} + \binom{89}{2n} = \sum_{n=0} ^{89} \binom{89}{n} = 2^{89}$, which also gives the correct answer. $\blacksquare$

 





\textbf{Solution 9}

We begin by computing
 \[\left(\prod_{k=1}^{90}\sin{k^{\circ}}\right)^2=\prod_{k=1}^{179}\sin{k^{\circ}}\]Let $z_{1}=e^{i1^{\circ}}$. Using the Euler's Formula $\sin{\theta}=\frac{e^{i\theta}-e^{-i\theta}}{2i}$, it follows that\begin{align*}\prod_{k=1}^{179}\sin{k^{\circ}}&=\frac{(z_{1}-\frac{1}{z_{1}})(z_{1}^{2}-\frac{1}{z_{1}^{2}})\ldots (z_{1}^{179}-\frac{1}{z_{1}^{179}})}{-2^{179}i}\\&=\frac{(1-z_{1}^{2})(1-z_{1}^{4})\ldots (1-z_{1}^{358})}{-(-1)^{179}\cdot z_{1}^{90\cdot 179}2^{179}i}\\&=\frac{(1-z_{1}^{2})(1-z_{1}^{4})\ldots (1-z_{1}^{358})}{2^{179}}\end{align*}
Let $f(x)=\frac{x^{180}-1}{x-1}$. Factoring by the Fundamental Theorem of Algebra and roots of unity, we find
 \[f(x)=(x-z_{1}^{2})(x-z_{1}^{4})\ldots (x-z_{1}^{358})=1+x+x^2+\ldots +x^{179}\]Thus,
 \[\prod_{k=1}^{179}\sin{k^{\circ}}=\frac{f(1)}{2^{179}}=\frac{180}{2^{179}}\]We proceed by computing
 \[\left(\prod_{k=1}^{45}\sin{2k^{\circ}}\right)^{2}=\prod_{k=1}^{89}\sin{2k^{\circ}}\]Let $z_{2}=e^{i2^{\circ}}$. We have\begin{align*}\prod_{k=1}^{89}\sin{2k^{\circ}}&=\frac{(z_{2}-\frac{1}{z_{2}})(z_{2}^{2}-\frac{1}{z_{2}^{2}})\ldots (z_{1}^{89}-\frac{1}{z_{1}^{89}})}{2^{89}i}\\&=\frac{(1-z_{1}^{2})(1-z_{1}^{4})\ldots (1-z_{1}^{178})}{(-1)^{89}\cdot z_{1}^{45\cdot 89}2^{89}i}\\&=\frac{(1-z_{1}^{2})(1-z_{1}^{4})\ldots (1-z_{1}^{178})}{2^{89}}\end{align*}
Let $g(x)=\frac{x^{90}-1}{x-1}$, we have
 \[g(x)=(x-z_{2}^{2})(x-z_{2}^{4})\ldots (x-z_{2}^{178})=1+x+x^2+\ldots +x^{89}\]Hence,
 \[\prod_{k=1}^{89}\sin{2k^{\circ}}=\frac{g(1)}{2^{89}}=\frac{90}{2^{89}}\]Therefore, the desired product evaluates to
 \[\prod_{k=1}^{45} \csc^2(2k-1)^\circ=\frac{\left(\prod_{k=1}^{45}\sin{2k^{\circ}}\right)^{2}}{\left(\prod_{k=1}^{90}\sin{k^{\circ}}\right)^2}=\frac{90}{2^{89}}\cdot \frac{2^{179}}{180}=2^{89}\]The requested sum is $2+89=\boxed{091}$

 




	\item (\textit{AIME I 2015}) For each integer $n \ge 2$, let $A(n)$ be the area of the region in the coordinate plane defined by the inequalities $1\le x \le n$ and $0\le y \le x \left\lfloor \sqrt x \right\rfloor$, where $\left\lfloor \sqrt x \right\rfloor$ is the greatest integer not exceeding $\sqrt x$. Find the number of values of $n$ with $2\le n \le 1000$ for which $A(n)$ is an integer.


 



 
 \textbf{Answer: }483



\textbf{Solution 1}

Let $n\ge 2$ and define $a(n) = \left\lfloor \sqrt n \right\rfloor$. For $2\le n \le 1000$, we have $1\le a(n)\le 31$. 

 For $a^2 \le x < (a+1)^2$ we have $y=ax$. Thus $A(n+1)-A(n)=a(n+\frac 12) = \Delta_n$ (say), and $\Delta_n$ is an integer if $a$ is even; otherwise $\Delta_n$ is an integer plus $\frac 12$.

 If $a=1$, $n\in \{1,2,3\}$ and $\Delta_n$ is of the form $k(n)+\frac 12$ so $A(n)$ is an integer when $n$ is even. 

 If $a=2$, $n\in\{4,\ldots , 8\}$ and $\Delta_n$ is an integer for all $n$. Since $A(3)$ is not an integer, so $A(n)$ is not an integer for any $n$. 

 If $a=3$, $n\in\{9,\ldots , 15\}$ and $\Delta_n$ is of the form $k(n)+\frac 12$. Since $A(8)$ is of the form $k+\frac 12$ so $A(n)$ is an integer only when $n$ is odd.

 If $a=4$, $n\in\{16,\ldots , 24\}$ and $\Delta_n$ is an integer for all $n$. Since $A(15)$ is an integer so $A(n)$ is an integer for all $n$. 

 Now we are back to where we started; i.e., the case $a=5$ will be the same as $a=1$ and so on. Thus, 
 \begin{align} a(n)\equiv 1\pmod 4 \qquad &\Longrightarrow \qquad A(n) \in \mathbb{Z} \textrm{ for even } n, \\ a(n)\equiv 2\pmod 4 \qquad &\Longrightarrow \qquad A(n) \not\in \mathbb{Z} \textrm{  for any } n, \\ a(n)\equiv 3\pmod 4 \qquad &\Longrightarrow \qquad A(n) \in \mathbb{Z} \textrm{ for odd  } n, \\ a(n)\equiv 0\pmod 4 \qquad &\Longrightarrow \qquad A(n) \in \mathbb{Z} \textrm{ for all } n. \end{align}

 For each $a$ there are $2a+1$ corresponding values of $n$: i.e., $n\in \{a^2, \ldots , (a+1)^2-1\}$.

 Thus, the number of values of $n$ corresponding to $(4)$ (i.e., $a(n)\equiv 0\pmod 4$) is given by 
 \[\sum_{\substack{a=4k \\ a\le 31}}(2a+1) = \sum_{k=1}^7 (8k+1)=231.\]

 The cases $(1)$ and $(3)$ combine to account for half the values of $n$ corresponding to odd values of $a(n)$; i.e.,
 \[\frac 12 \cdot \sum_{\substack{a=2k+1 \\ a\le 31}} (2a+1) = \sum_{k=0}^{15} (2k+\frac 32) = 264\]However, this also includes the odd integers in $\{1001, \ldots , 1023\}$. Subtracting $12$ to account for these, we get the  number of values of $n$ corresponding to cases $(1)$ and $(3)$ to be $264-12=252$.

 Adding the contributions from all cases we get our answer to be $231+252= \boxed{483}$.

 





\textbf{Solution 2}

By considering the graph of this function, it is shown that the graph is composed of trapezoids ranging from $a^2$ to $(a+1)^2$ with the top made of diagonal line $y=ax$. The width of each trapezoid is $3, 5, 7$, etc. Whenever $a$ is odd, the value of $A(n)$ increases by an integer value, plus $\frac{1}{2}$. Whenever $a$ is even, the value of $A(n)$ increases by an integer value. Since each trapezoid always has an odd width, every value of $n$ is not an integer when $a \pmod{4} \equiv 2$, and is an integer when $a \pmod{4} \equiv 0$. Every other value is an integer when $a$ is odd. Therefore, it is simply a matter of determining the number of values of $n$ where $a \pmod{4} \equiv 0$ ($(5^2-4^2)+(9^2-8^2)+...+(29^2-28^2)$), and adding the number of values of $n$ where $a$ is odd ($\frac{(2^2-1^2)+(4^2-3^2)+...+(30^2-29^2)+(1000-31^2)}{2}$). Adding the two values gives $231+252=\boxed{483}$.

 \textbf{"Step" Solution} First, draw a graph of the function. Note that it is just a bunch of line segments. Since we only need to know whether or not $A(n)$ is an integer, we can take the area of each piece from some $x$ to $x+1$ (mod 1), aka the piece from $2$ to $3$ has area $\frac{1}{2} (\mod 1)$. There are some patterns. Every time we increase $n$ starting with $2$, we either add $0 (\mod 1)$ or $\frac{1}{2} (\mod 1)$. We look at $\lfloor \sqrt{x} \rfloor$ for inspiration. Every time this floor (which is really the slope) is odd, there is always an addition of $\frac{1}{2} (\mod 1)$, and whenever that slope is even, that addition is zero.

 Take a few cases. For slope $=1$, we see that only one value satisfies. Because the last value, $n=4$, fails, and the numbers $n$ which have a slope of an even number don't change this modulus, all these do not satisfy the criterion. The pattern then comes back to the odds, and this time $\lfloor \frac{7}{2} \rfloor + 1 = 4$ values work. Since the work/fail pattern alternates, all the $n$s with even slope, $[17, 25]$, satisfy the criterion. This pattern is cyclic over period 4 of slopes.

 Even summation of working cases: $9+17+25+...+57 = 231$.Odd summation: $1+4+5+8+9+12...+29$ and plus the $20$ cases from $n=[962, 1000]$: $252$. Answer is $\boxed{483}$.

 





\textbf{Solution 4 (Integration)}

We see that
 \[A(n)=\int_{1}^{n}x\lfloor \sqrt{x}\rfloor \,dx\]Let $n\in [a^2, (a+1)^2)$ for $1\le a \le 31$, we see that\begin{align*}\int_{1}^{n}x\lfloor \sqrt{x}\rfloor\,dx&= \left(\sum_{k=1}^{a-1}\int_{k^2}^{(k+1)^2}kx \,dx\right)+\int_{a^2}^{n}ax \,dx\\&=\left(\sum_{k=1}^{a-1}\frac{k(2k+1)\bigl[(k+1)^2+k^2\bigr]}{2}\right)+\frac{a(n^2-a^4)}{2}\end{align*}
Let $S(a)= \sum_{k=1}^{a-1}\frac{k(2k+1)\bigl[(k+1)^2+k^2\bigr]}{2}$. Let $\{x\}=x-\lfloor x\rfloor$ denote the fractional part of $x$. We see that\begin{align*}&2\mid k \implies \frac{k(2k+1)[(k+1)^2+k^2]}{2}\in \mathbf{Z}\\&2\nmid k \implies \{\frac{k(2k+1)[(k+1)^2+k^2]}{2}\}=\frac{1}{2}\end{align*}
By casework on $a\bmod{4}$, we find:\begin{array}{c|c|c|c}a \bmod 4 & \{S(a)\} & \text{desired }\{\frac{a(n^2-a^4)}{2}\} & \text{condition needed for } n \\ \hline0 & 0 & 0 & \text{all }n \\1 & 0 & 0 & \text{odd} \\2 & 1/2 & 1/2 & \text{no solution} \\3 & 1/2 & 1/2 & \text{even}\end{array}There are strictly $2a+1$ integers in the interval $[a^2,(a+1)^2)$. For \textbf{Case 1}, the number of values of $n$ is
 \[\sum_{k=1}^{7}8k+1=231\]For an odd integer $a$, the number of odd integers in the interval $[a^2,(a+1)^2)$ is $a+1$. For \textbf{Case 2}, we have
 \[\sum_{k=0}^{7}4k+2=128\]There are no desired values of $n$ for \textbf{Case 3}. 

 For an odd integer $a$, the number of even integers in the interval $[a^2,(a+1)^2)$ is $a$. For \textbf{Case 4}, we have
 \[\sum_{k=0}^{7}4k+3=136\]However, we shall exclude $n=1$ and $n\in\{1002, 1004, \ldots, 1022\}$. The total number of values of $n$ is
 \[231+128+136-1-11=\boxed{483}\]

 


 


 




	\item (\textit{AIME I 2015}) A block of wood has the shape of a right circular cylinder with radius $6$ and height $8$, and its entire surface has been painted blue.  Points $A$ and $B$ are chosen on the edge of one of the circular faces of the cylinder so that $\overarc{AB}$ on that face measures $120^\text{o}$.  The block is then sliced in half along the plane that passes through point $A$, point $B$, and the center of the cylinder, revealing a flat, unpainted face on each half.  The area of one of these unpainted faces is $a\cdot\pi + b\sqrt{c}$, where $a$, $b$, and $c$ are integers and $c$ is not divisible by the square of any prime.  Find $a+b+c$.


 
                    \begin{figure}[H]
                        \centering
                        \includegraphics[width=0.5\textwidth]{Problem/AIME_2015_I_15_Problem_Figure_1.png}
                    \end{figure}
                    


 
 \textbf{Answer: }053



\textbf{Solution 1}


                    \begin{figure}[H]
                        \centering
                        \includegraphics[width=0.5\textwidth]{Solution/AIME_2015_I_15_Solution_Figure_1.png}
                    \end{figure}
                    Label the points where the plane intersects the top face of the cylinder as $C$ and $D$, and the center of the cylinder as $O$, such that $C,O,$ and $A$ are collinear. Let $N$ be the center of the bottom face, and $M$ the midpoint of $\overline{AB}$. Then $ON=4$, $MN=3$ (because of the 120 degree angle), and so $OM=5$. 

 Project $C$ and $D$ onto the bottom face to get $X$ and $Y$, respectively. Then the section $ABCD$ (whose area we need to find), is a stretching of the section $ABXY$ on the bottom face. The ratio of stretching is $\frac{OM}{MN}=\frac{5}{3}$, and we do not square this value when finding the area because it is only stretching in one direction*. Using 30-60-90 triangles and circular sectors, we find that the area of the section $ABXY$ is $18\sqrt{3}\ + 12 \pi$. Thus, the area of section $ABCD$ is $20\pi + 30\sqrt{3}$, and so our answer is $20+30+3=\boxed{053}$.

 \textbf{*Easier way to think of scale} Notice that we only scale one sector, particularly $CD = AB$ by $1$ from $BC'$ and $AD'$ (or $AY$ and $BX$). The scale between the remaining curvatures are of the factor $\frac{5}{3}$. The area of the smaller figure is $18\sqrt{3} + 12\cdot \pi$, and the area of the desired figure is the area of the smaller figure times both scale factors, or $(18\sqrt{3} + 12\cdot \pi) \cdot (\frac{5}{3} \cdot 1)$, giving the requested result.

 





\textbf{Solution 2}

Label the points same as in the first sentence above. Consider a view of the cylinder such that height is disregarded, i.e. a top view. From this view, note that Cylinder $O$ has become a circle with $\overarc{AB}$ = $\overarc{CD}$ = $120^\text{o}$. Using 30-60-90 triangles, we get rectangle $ABCD$ to have a horizontal component of $6$. Now, consider a side view, such that $A$ and $B$ coincide at the bottom of the diagram. From this view, consider the right triangle composed of hypotenuse $AD$ and a point along the base of the viewpoint, which will be labeled as $E$. From the top view, $AE = 6$. Because of the height of the cylinder, $DE$ is equal to $8$. This makes $AD$ equal to $10$. 

 
Now, the use of simple calculus is required. Conceptualize an infinite number of lines perpendicular to $AE$ intersecting both $AE$ and $AD$. Consider the area between point $A$ and the first vertical line. Label the point where the line intersects AE as E', and the point where the line intersects AD as D'. The area of the part of the initial unpainted face within these two positions approaches a rectangle with length AD' and width $w$. The area of the base within these two positions approaches a rectangle with length AE' and width $w$. The ratio of AD':AE' is 10:6, since the ratio of AD:AE is 10:6. This means that the area of the initial unpainted surface within these two positions to the area of the base within these two positions is equal to 10$w$:6$w$ = 10:6. Through a similar argument, the areas between each set of vertical lines also maintains a ratio of 10:6. Therefore, the ratio of the area we wish to find to the area of the base between AB and CD (from the top perspective) is 10:6. Using 30-60-90 triangles and partial circles, the area of the base between AB and CD is calculated to be $18\sqrt{3}\ + 12 \pi$. The area of the unpainted surface therefore becomes $20\pi + 30\sqrt{3}$, and so our answer is $\boxed{053}$.

 





\textbf{Solution 3}

This problem can be calculus-bashed (for those like me who never noticed the surface was merely a stretch of its projection). Label points as in the first paragraph of Solution 1 ($A$ and $B$ as given, $M$ the midpoint of $AB$, $O$ the center of the cylinder, $T$ the center of the bottom face of the cylinder). Because of the 120 degrees and right triangle calculations, we find $MT$ = 3, $OT$ = 4, $OM$ = 5). We will be integrating with respect to the y-coordinate which we define as distance downwards from $O$ (in this system, the $y$-coordinate of the bottom face would be 4). 

 We note that by similar triangles, we have that the length from $O$ to the point on the unpainted surface of coordinate $y$ is $\ell = \frac{5}{4} y$, and therefore $d\ell = \frac{5}{4} dy$. Define the segment $A'B'$ to be the intersection of the painted surface with the circular cross section of the cylinder of coordinate $y$, with endpoints $A'$ and $B'$ and midpoint $M'$, with $T'$ the center of this circular cross section. Then, by similar triangles, $T'M' = \frac{3}{4} y$ and thus 
 \[A'B' = 2A' M' = 2 \sqrt{ 6^2 - \left( \frac{3}{4}y\right)^2 } = \frac{3}{2} \sqrt{ 64 - y^2 }\]. We know that $A'B'$ is perpendicular to $OM$. 

 Now we can set up our integral: we will integrate $y$ from 0 to 4 and multiply by two because the total height is 8.
 \[A = 2\int_0^4 \left(\frac{3}{2}\sqrt{ 64 - y^2 }\right) \left( \frac{5}{4} dy\right)\]
 \[A = \frac{15}{4} \int_0^4 \sqrt{ 64 - y^2 }dy\]

 Then we substitute $y = 8\sin \theta$ with $dy = 8 \cos \theta d \theta$, changing the bounds to 0 to $\frac{\pi}{6}$ as appropriate.
 \[A = \frac{15}{4} \int_0^\frac{\pi}{6} \sqrt{ 64 - 64\sin^2 \theta }\cdot 8\cos\theta d\theta\]
 \[A = 240 \int_0^\frac{\pi}{6} \cos^2 \theta d\theta\]
 \[A = 240 \left[ \frac{\theta}{2} + \frac{\sin 2\theta}{4} \right]_0^\frac{\pi}{6}  = 240 \left[ \frac{\pi}{12} + \frac{\sqrt{3}}{8} \right] = 20{\pi} + 30\sqrt{3}\]

 Therefore, $a + b + c = 20 + 30 + 3 = \boxed{053}$.

 


 





\textbf{Solution 4 - Author : Shiva Kumar Kannan - In Progress}

Extend the cylinder such that the cylinder has a height of $16$ and the same radius of $6$. Let $Q$ be a point on the circumference of the top face of the cylinder. Let $R$ be a point on the bottom face of the cylinder, and the farthest point from A on the cylinder. Let a plane $P$ pass through points $Q$ and $R$. The cross section formed by the intersection of the plane and this stretched cylinder, is an ellipse. Notice that the cross section area we want to find is a part the part of this ellipse contained in the original cylinder. 

 
The semi-major axis length of the ellipse is 
 \[\sqrt{6^2 + 8^2} = 10\]. The semiminor axis length is just the radius of the cylinder, which is $6$. 

 
If we take the plane of this ellipse to be the $2$ dimensional Cartesian plane and take the center of this ellipse as $(0, 0)$, this ellipse has the equation : 
 \[\frac{x^2}{10^2} + \frac{y^2}{6^2} = 1\].

 
Notice that the cross section area we want to find is the area bounded by this ellipse from $x = -5$ to $x = 5$. Thus we must integrate the equation of the ellipse from $-5$ to $5$. That will give us the positive area bounded by the ellipse and the $x$ axis. We want to double it, as we want both the positive and negative area as that is our cross section. 

 
Rewrite the ellispe equation as 
 \[y = \frac{3}{5} \sqrt{100 - x^2}\]. We want to find 
 \[2\int_{-5}^{5} \left(\frac{3}{5}\sqrt{100 - x^2 }\right) \left( dx\right)\].

 
Evaluate this integral : It comes out to be 
 \[20\pi + 30\sqrt{3}\]. Finally, the answer we give, is 
 \[20 + 30 + 3 = \boxed{053}\].

 Motivation : Notice that a cross section of any diagonal plane not going through the top or bottom face of the cylinder is an ellipse. When the plane does cross the bottom and top faces, the resulting cross section will be a portion of an ellipse. When the plane is perpendicular to its top and bottom faces, the resulting cross section is a rectangle. So, from the moment the cross section starts cutting both circular faces, to the moment it is perpendicular to the circular bases, it goes from a full ellipse, to a rectangle. Hence, the cross section we want is a part of the ellipse.

 





\textbf{Solution 5}

Let $O$ be the center of the circle where $A,B$ lie, and let $\mathcal{P}$ be the plane intersecting the cylinder. Also, let $O'=(0,0,4)$. We establish a 3D coordinate system with $O$ being the origin. Here is our setup:\begin{align*}A&=(-3,3\sqrt{3},0)\\B&=(-3,-3\sqrt{3},0)\\O&=(0,0,0)\\O'&=(0,0,4)\\\text{Cylinder}&: x^2+y^2=36, 0\le z\le 8\\\mathcal{P}&: Ax+By+Cz+D=0\end{align*}
Points $A,B,O'$ lie on the plane $\mathcal{P}$. We obtain the following system:\begin{align*}4C+D&=0\\-3A+3\sqrt{3}B+D&=0\\-3A-3\sqrt{3}B+D&=0\\\end{align*}
We see that $B=0, D=-4C, D=3A$. WLOG, let $C=-3$. It follows that the equation for $\mathcal{P}$ is
 \[4x-3z+12=0\]Or
 \[z=\frac{4}{3}x+4\]The area element $dS$ is calculated as follows:
 \[dS=\sqrt{1+\left(\frac{\partial z}{\partial x}\right)^2+\left(\frac{\partial z}{\partial y}\right)^2}\,dx\,dy=\sqrt{1+\left(\frac{4}{3}\right)^2+0^2}\,dx\,dy=\frac{5}{3}\,dx\,dy\]Recall the formula to calculate the cross-section:
 \[\iint_{D}\sqrt{1+\left(\frac{\partial z}{\partial x}\right)^2+\left(\frac{\partial z}{\partial y}\right)^2}\,dx\,dy=\iint_{D}\frac{5}{3}\,dx\,dy\]We find $x=\frac{3}{4}z-3$, and $0\le z\le 8\implies -3\le x\le 3$. From the equation of the cylinder, $y=\pm\sqrt{36-x^2}$, we have
 \[D=\{(x,y)\mid -3\le x\le 3,\ -\sqrt{36-x^2}\le y\le \sqrt{36-x^2}\}\]Thus, the area of the cross section is
 \[\int_{-3}^{3}\int_{-\sqrt{36-x^2}}^{\sqrt{36-x^2}}\frac{5}{3}\,dy\,dx=\frac{10}{3}\int_{-3}^{3}\sqrt{36-x^2}\,dx\]The former integral represents the shaded area in the following semicircle with radius $6$:
                    \begin{figure}[H]
                        \centering
                        \includegraphics[width=0.5\textwidth]{Solution/AIME_2015_I_15_Solution_Figure_1.png}
                    \end{figure}
                    Adding the $\frac{1}{6}$ sector and 2 special triangles gives the area
 \[\frac{1}{6}\cdot 36\pi+2\cdot \frac{1}{2}\cdot 3\cdot 3\sqrt{3}=6\pi+9\sqrt{3}\]The desired area is hence
 \[\frac{10}{3}\int_{-3}^{3}\sqrt{36-x^2}\,dx=\frac{10}{3}(6\pi+9\sqrt{3})=20\pi+30\sqrt{3}\]The requested sum is $20+30+3=\boxed{053}$.

 \textbf{Remark} Note that you can derive the formula for $dS$ as follows:

 Suppose the surface is given by \(z=f(x,y)\). At a point \((x,y,f(x,y))\), we take small increments \(dx\) and \(dy\) in the \(x\)- and \(y\)-directions, respectively.

 When \(y\) is fixed and \(x\) increases by \(dx\), the displacement vector on the surface along the \(x\)-direction is
 \[\vec r_x=\left(dx,\ 0,\ \frac{\partial f}{\partial x}\,dx\right),\]because the change in \(z\) is approximately \(\frac{\partial f}{\partial x}dx\).

 Similarly, when \(x\) is fixed and \(y\) increases by \(dy\), the displacement along the \(y\)-direction is
 \[\vec r_y=\left(0,\ dy,\ \frac{\partial f}{\partial y}\,dy\right).\]

 These two vectors span an infinitesimal parallelogram lying in the tangent plane to the surface; its area is precisely the surface element \(dS\). The area of this parallelogram is equal to the magnitude of the cross product of the two side vectors.

 Compute the cross product:\begin{vmatrix}\mathbf i & \mathbf j & \mathbf k\\dx & 0 & \frac{\partial f}{\partial x}dx\\0 & dy & \frac{\partial f}{\partial y}dy\end{vmatrix}Computing yields $\vec r_x\times \vec r_y=\left(-\frac{\partial f}{\partial x},\ -\frac{\partial f}{\partial y},\ 1\right)dx\,dy$

 Taking the magnitude gives
 \[|\vec r_x\times \vec r_y|=\sqrt{1+\left(\frac{\partial f}{\partial x}\right)^2+\left(\frac{\partial f}{\partial y}\right)^2}\,dx\,dy.\]

 Therefore, the surface area element is
 \[\boxed{dS=\sqrt{1+\left(\frac{\partial z}{\partial x}\right)^2+\left(\frac{\partial z}{\partial y}\right)^2}\,dx\,dy}.\]

 \href{https://mathproblemsolvingskills.wordpress.com/}{MathProblemSolvingSkills.com}

 


 




\end{enumerate}